\documentclass[lettersize,journal]{IEEEtran} 
\usepackage[colorlinks=true]{hyperref} 
% \usepackage{orcidlink} 
\usepackage[english]{babel} 
\usepackage{graphicx} 
% \graphicspath{{ {./images/} }} 
\usepackage{amsmath,amsfonts}
\usepackage{amssymb} 
\usepackage{subcaption} 
\usepackage{cite}
\usepackage{multirow}
\usepackage{enumitem}
\usepackage{mathtools}
% \usepackage[compact]{titlesec} 
% \usepackage{ragged2e}
% \AtBeginDocument{
% \setlength\abovedisplayskip{3pt}  
% \setlength\belowdisplayskip{3pt}} 
% \setlength{\textfloatsep}{1pt} 
% \usepackage[font=small]{caption}
% \captionsetup[figure]{labelfont=small,name={Fig.},labelsep=period}

\begin{document}

\title{Image-based Visual Servoing Scheme for Soft Precise Landing on Mobile Platform}

\author{Shubham~Singhal}

% \markboth{Data-driven learning}%
% {}

\maketitle
\section*{Nomenclature}

\begin{tabular}{@{}l @{\quad=\quad} l@{}}
$t$ & time, s \\
$m$ & mass, kg \\
$\boldsymbol{r}$ & radial position vector, m \\
$\boldsymbol{v}$ & linear velocity vector, m/s \\
$\boldsymbol{\omega}$ & angular velocity vector, rad/s \\
$\boldsymbol{F}$ & force vector, N \\
$J$ & moment of inertia matrix, kg m$^2$ \\
$\boldsymbol{\tau}$ & torque vector, N m \\
$T$ & total thrust of the rotors, N \\
$\hat{\boldsymbol{r}}$ & image point \\
$\boldsymbol{s}$ & normalized position vector, rad \\
$\alpha$ & image orientation, rad \\
$\boldsymbol{h}$ & optical flow, rad/s \\
$\boldsymbol{q}$ & system output vector \\
$\boldsymbol{u}$ & system input vector \\
$\beta$ & uncertainty parameter \\
$\boldsymbol{c}$ & vector representing known system dynamics \\
$\boldsymbol{d}$ & vector representing unknown system dynamics \\
$\boldsymbol{p}$ & performance function vector \\
$V$ & Lyapunov function
\end{tabular}

\section{PRELIMINARIES \& PROBLEM FORMULATION} \label{background: section}
This section outlines the system dynamics and control strategy of an underactuated quadrotor UAV platform using scale-independent target parameters and formulates the problem to address the challenges in vision-based precise landing.

\subsection{Dynamics of Quadrotor-Target System} \label{quadrotor dynamics: section}
Achieving soft precise landing on a moving target requires a thorough understanding of the system dynamics and control architecture of an underactuated quadrotor UAV platform. 
First, consider a body-fixed frame $\mathcal{B}=\{O_{\text{b}} ~X_{\text{b}} ~Y_{\text{b}} ~Z_{\text{b}}\}$ rigidly attached to the UAV in forward-right-down (FRD) configuration.
Let the linear and angular velocities of the UAV be denoted as $\,^\mathcal{B} \boldsymbol{v}_{\text{b}}$ and $\,^\mathcal{B} \boldsymbol{\omega}_{\text{b}}$ respectively.
Now, using the Newton-Euler formulation, the UAV flight dynamics can be expressed as follows:
\begin{align}
\label{quadrotor dynamics: equation}
    m \,^\mathcal{B} \dot{\boldsymbol{v}}_{\text{b}} &= \,^\mathcal{B} \boldsymbol{F_{\text{ext}}} - \,^\mathcal{B} \boldsymbol{\omega}_{\text{b}} \times (m \,^\mathcal{B} \boldsymbol{v}_{\text{b}})  \nonumber \\
    J \,^\mathcal{B} \dot{\boldsymbol{\omega}}_{\text{b}} &= \,^\mathcal{B} \boldsymbol{\tau_{\text{ext}}} - \,^\mathcal{B} \boldsymbol{\omega}_{\text{b}} \times (J \,^\mathcal{B} \boldsymbol{\omega}_{\text{b}}),
\end{align}
where $m$ is the mass,  and $J = \text{diag}(J_{xx}~J_{yy}~J_{zz})$ is the inertia matrix.
The external force $\,^\mathcal{B} \boldsymbol{F_{\text{ext}}}$ includes gravitational force $\,^\mathcal{B} \boldsymbol{F_{\text{g}}}$, the total body thrust $\,^\mathcal{B} \boldsymbol{F_{\text{t}}} = \begin{bmatrix} 0 & 0 & \,^\mathcal{B} T_u\end{bmatrix}^\top $ with $\,^\mathcal{B} T_u$ being the total rotor control thrust, and the unknown model uncertainties and disturbances $\,^\mathcal{B} \boldsymbol{F_{\text{d}}}$ such that $\,^\mathcal{B} \boldsymbol{F_{\text{ext}}} = \,^\mathcal{B} \boldsymbol{F_{\text{g}}} + \,^\mathcal{B} \boldsymbol{F_{\text{t}}} + \,^\mathcal{B} \boldsymbol{F_{\text{d}}}$.
Similarly, the external moment $\,^\mathcal{B} \boldsymbol{\tau_{\text{ext}}}$ includes the net control moments of the four rotors $\,^\mathcal{B} \boldsymbol{\tau}_u$, and the unknown model uncertainties and disturbances $\,^\mathcal{B} \boldsymbol{\tau_{\text{d}}}$ such that $\,^\mathcal{B} \boldsymbol{\tau_{\text{ext}}} = \,^\mathcal{B} \boldsymbol{\tau}_u + \,^\mathcal{B} \boldsymbol{\tau_{\text{d}}}$.
Next, consider a target located at point $\,^\mathcal{B} \boldsymbol{r}_{\text{t}}$ with respect to the origin $O_{\text{b}}$ in body frame $\mathcal{B}$. It moves with a relative velocity $\,^\mathcal{B} \boldsymbol{v}_{\text{t/b}}$ and relative acceleration $\,^\mathcal{B} \boldsymbol{a}_{\text{t/b}}$. Then, the target kinematics in frame $\mathcal{B}$ is obtained as follows:
\begin{align} \label{target kinematics: equation}
    \,^\mathcal{B} \dot{\boldsymbol{r}}_{\text{t}} &= \,^\mathcal{B} \boldsymbol{v}_{\text{t/b}} - \,^\mathcal{B} \boldsymbol{\omega}_{\text{b}} \times \,^\mathcal{B} \boldsymbol{r}_{\text{t}}, \nonumber \\
    \,^\mathcal{B} \dot{\boldsymbol{v}}_{\text{t}} &= \,^\mathcal{B} \boldsymbol{a}_{\text{t/b}} - \,^\mathcal{B} \dot{\boldsymbol{\omega}}_{\text{b}} \times \,^\mathcal{B} \boldsymbol{r}_{\text{t}} - \,^\mathcal{B} \boldsymbol{\omega}_{\text{b}} \times \left( \,^\mathcal{B} \boldsymbol{\omega}_{\text{b}} \times \,^\mathcal{B} \boldsymbol{r}_{\text{t}} \right) \nonumber \\
    &~~~ - 2 \left( \,^\mathcal{B} \boldsymbol{\omega}_{\text{b}} \times \,^\mathcal{B} \boldsymbol{v}_{\text{t/b}} \right)
\end{align}

\subsection{Pose-based Flight Controller}
In order to achieve target tracking using a pose-based flight controller, the overall control architecture is a cascaded PID controller with outer-loop operating in the inertial frame $\mathcal{I}=\{O_{\text{i}} ~X_{\text{i}} ~Y_{\text{i}} ~Z_{\text{i}}\}$ and the inner-loop in the body frame $\mathcal{B}$.
The outer-loop controller is designed to generate desired velocity $\,^\mathcal{I}\boldsymbol{v}_{\text{d}}$ command that drives the UAV from the instantaneous position $\,^\mathcal{I} \boldsymbol{r}_{\text{b}}$ to the known target position $\,^\mathcal{I}\boldsymbol{r}_{\text{t}}$.
Then, it generates desired acceleration $\,^\mathcal{I}\boldsymbol{a}_{\text{d}}$ command that drives the instantaneous UAV velocity $\,^\mathcal{I} \boldsymbol{v}_{\text{b}}$ to the desired velocity $\,^\mathcal{I}\boldsymbol{v}_{\text{d}}$.
The control laws are shown below:
\begin{align} \label{outer_loop_control: equation}
    \,^\mathcal{I}\boldsymbol{v}_{\text{d}} =& - K_{rp} \,^\mathcal{I}\boldsymbol{r}_{\text{e}} - K_{ri} \int \,^\mathcal{I}\boldsymbol{r}_{\text{e}} dt - K_{rd} \frac{d\,^\mathcal{I}\boldsymbol{r}_{\text{e}}}{dt} \nonumber \\
    \,^\mathcal{I}\boldsymbol{a}_{\text{d}} =&  K_{vp} \,^\mathcal{I}\boldsymbol{v}_{\text{e}} - K_{vi} \int \,^\mathcal{I}\boldsymbol{v}_{\text{e}} dt - K_{vd} \frac{d\,^\mathcal{I}\boldsymbol{v}_{\text{e}}}{dt},
\end{align}
where $\,^\mathcal{I}\boldsymbol{r}_{\text{e}} = \,^\mathcal{I}\boldsymbol{r}_{\text{b}} - \,^\mathcal{I}\boldsymbol{r}_{\text{t}}$ represents the relative position between the UAV and the target and $\,^\mathcal{I}\boldsymbol{v}_{\text{e}} = \,^\mathcal{I}\boldsymbol{v}_{\text{b}} - \,^\mathcal{I}\boldsymbol{v}_{\text{d}}$.

Since the centrifugal force $\,^\mathcal{B} \boldsymbol{\omega}_{\text{b}} \times (m \,^\mathcal{B} \boldsymbol{v}_{\text{b}})$, given in \eqref{quadrotor dynamics: equation}, is nullified in the inertia frame, the desired acceleration $\,^\mathcal{I}\boldsymbol{a}_{\text{d}}$ is assumed to be achieved solely by the total rotor control thrust $\,^\mathcal{B} T_u$ and its direction, after compensating for gravity and disturbances.
The relationship between the desired acceleration ($\,^\mathcal{I}\boldsymbol{a}_{\text{d}}$) and the thrust force applied along the desired body Z-axis ($\,^\mathcal{B} T_u$) is derived as:
\begin{align} \label{a_des: equation}
    \,^\mathcal{I} \boldsymbol{a}_{\text{d}} &=  \,^\mathcal{I}R_{\mathcal{B}} \,^\mathcal{B} \boldsymbol{F}_{\text{t}}/m = \,^\mathcal{I}R_{\mathcal{B}} \begin{bmatrix} 0 & 0 & \,^\mathcal{B} T_u /m\end{bmatrix}^\top \nonumber \\
    &= \frac{\,^\mathcal{B} T_u}{m}\begin{bmatrix} S_\psi S_{\theta} C_{\phi} -C_\psi S_{\phi} \\ C_\psi S_{\theta} C_{\phi} +S_\psi S_{\phi} \\- C_{\theta} C_{\phi} \end{bmatrix},
\end{align}
where the rotation matrix $^\mathcal{I}R_\mathcal{B}$ transforms vectors from body frame ($\mathcal{B}, FRD$) to the inertial frame ($\mathcal{I}, ENU$), is given as.
\begin{align}
\label{I_R_B: equation}
    ^\mathcal{I}R_\mathcal{B} = & ^{ENU}R_{NED} R_z(\psi) R_y(\theta) R_x(\phi) \nonumber \\
    = & \begin{bmatrix} S_\psi C_\theta &S_\psi S_\theta S_\phi +C_\psi C_\phi &S_\psi S_\theta C_\phi -C_\psi S_\phi \\C_\psi C_\theta &C_\psi S_\theta S_\phi -S_\psi C_\phi &C_\psi S_\theta C_\phi +S_\psi S_\phi \\ S_\theta &- C_\theta S_\phi & - C_\theta C_\phi \end{bmatrix},
\end{align}
in which $S_x = sin(x)$ and $C_x = cos(x)$.

The desired roll (${\phi}_{\text{d}}$), pitch (${\theta}_{\text{d}}$), and total thrust ($\,^\mathcal{B} T_u$) commands are then analytically derived by inverting the acceleration $\,^\mathcal{I}\boldsymbol{a}_{\text{d}}$ in \eqref{a_des: equation} using its components ($\,^\mathcal{I} {a}_{\text{d}_x}, \,^\mathcal{I} {a}_{\text{d}_y}, \,^\mathcal{I} {a}_{\text{d}_z}$):
\begin{align}
\label{euler_thrust_des: equation}
    {\theta}_{\text{d}} = & \tan^{-1}\left(\frac{- S_\psi \,^\mathcal{I} {a}_{\text{d}_x} - C_\psi \,^\mathcal{I} {a}_{\text{d}_y}}{\,^\mathcal{I} {a}_{\text{d}_z}} \right) \nonumber \\
    {\phi}_{\text{d}} = & \tan^{-1}\left(\frac{C_{\theta_{\text{d}}} \left(C_\psi \,^\mathcal{I} {a}_{\text{d}_x} - S_\psi \,^\mathcal{I} {a}_{\text{d}_y} \right)}{\,^\mathcal{I} {a}_{\text{d}_z}} \right) \nonumber \\
    \,^\mathcal{B} {T}_u = & - \frac{m \,^\mathcal{I} {a}_{\text{d}_z}}{C_{\theta_{\text{d}}} C_{\phi_{\text{d}}}},
\end{align}
Combining the derived attitude ($\phi_{\text{d}}, \theta_{\text{d}}$) with the known target heading $\psi_{\text{t}}$, the desired attitude $\boldsymbol{\Phi}_{\text{d}} = [\phi_{\text{d}}, \theta_{\text{d}}, \psi_{\text{t}}]^\top$ and thrust ($\,^\mathcal{B} T_u$) commands are then passed to the inner-loop controller to execute the required control effort.

The inner-loop controller uses a cascade of PID loops to regulate the attitude ($\boldsymbol{\Phi}_{\text{b}}$) and angular body rates ($\,^\mathcal{B} \boldsymbol{\omega}_{\text{b}}$) of the UAV.
The goal is to generate desired attitude rates ($\dot{\boldsymbol{\Phi}}_{\text{d}}$) command that helps the UAV to accurately follow the desired attitude ($\boldsymbol{\Phi}_{\text{d}}$) command generated by the outer-loop controller.
The desired attitude rates ($\dot{\boldsymbol{\Phi}}_{\text{d}}$) are transformed into the desired angular body rates ($\,^\mathcal{B} \boldsymbol{\omega}_{\text{d}}$) using the kinematic transformation matrix $W(\phi_{\text{d}} ,\theta_{\text{d}} )$:
\begin{equation}
\label{kinematic transformation: equation}
    \,^\mathcal{B} \boldsymbol{\omega}_{\text{d}} = W(\phi_{\text{d}} ,\theta_{\text{d}} ) \dot{\boldsymbol{\Phi}}_{\text{d}},
\end{equation}
where matrix $W(\phi_{\text{d}} ,\theta_{\text{d}} )$ relates the inertial Euler rates to the body frame angular rates, based on the desired roll ($\phi_{\text{d}}$) and pitch ($\theta_{\text{d}}$) angles as shown below:
\begin{equation}
    W(\phi_{\text{d}} ,\theta_{\text{d}} )=\begin{bmatrix} 0 & C_{\phi_{\text{d}}} & S_{\phi_{\text{d}}} C_{\theta_{\text{d}}} \\ 1 & 0 & -S_{\theta_{\text{d}}} \\ 0 & S_{\phi_{\text{d}}} & - C_{\phi_{\text{d}}} C_{\theta_{\text{d}}} \end{bmatrix}
\end{equation}

Then, the inner-loop controller generates desired angular acceleration $\,^\mathcal{B} \dot{\boldsymbol{\omega}}_{\text{d}}$ command that drives the current body rate $\,^\mathcal{B} \boldsymbol{\omega}_{\text{b}}$ of the UAV to the desired body rate $\,^\mathcal{B} \boldsymbol{\omega}_{\text{d}}$.
The control laws are shown below:
\begin{align}
    \dot{\boldsymbol{\Phi}}_{\text{d}} =& - K_{\Phi p} \boldsymbol{\Phi}_{\text{e}} - K_{\Phi i} \int \boldsymbol{\Phi}_{\text{e}} dt - K_{\Phi d} \frac{d\boldsymbol{\Phi}_{\text{e}}}{dt} \nonumber \\
    \,^\mathcal{B} \dot{\boldsymbol{\omega}}_{\text{d}} =& - K_{\omega p} \,^\mathcal{B} \boldsymbol{\omega}_{\text{e}} - K_{\omega i} \int \,^\mathcal{B} \boldsymbol{\omega}_{\text{e}} dt - K_{\omega d} \frac{d\,^\mathcal{B} \boldsymbol{\omega}_{\text{e}}}{dt},
\end{align}
where $\boldsymbol{\Phi}_{\text{e}} = \boldsymbol{\Phi}_{\text{b}} - \boldsymbol{\Phi}_{\text{d}}$ and $\boldsymbol{\omega}_{\text{e}} = \,^\mathcal{B} \boldsymbol{\omega}_{\text{b}} - \,^\mathcal{B} \boldsymbol{\omega}_{\text{d}}$.
The desired angular acceleration $\,^\mathcal{B} \dot{\boldsymbol{\omega}}_{\text{d}}$ is then used with the rigid body dynamics equation \eqref{quadrotor dynamics: equation} to calculate the necessary desired control torque ($\,^\mathcal{B} \boldsymbol{\tau}_u$):
\begin{equation}
\label{tau_des: equation}
    \,^\mathcal{B} \boldsymbol{\tau}_u = J \,^\mathcal{B} \dot{\boldsymbol{\omega}}_{\text{d}} + \,^\mathcal{B} \boldsymbol{\omega}_{\text{d}} \times (J \,^\mathcal{B} \boldsymbol{\omega}_{\text{d}})
\end{equation}
The final commands sent to the motor mixing algorithm are the combined thrust-torque command $\boldsymbol{u} = [\,^\mathcal{B} {T}_u, \,^\mathcal{B} \boldsymbol{\tau}_u^\top]^\top$.

\subsection{Virtual Feature Points}
Consider a camera frame $\mathcal{C}=\{O_{\text{c}} ~X_{\text{c}} ~Y_{\text{c}} ~Z_{\text{c}}\}$ rigidly attached to a monocular camera such that its $X_{\text{c}}$ and $Y_{\text{c}}$ axes are aligned with the image horizontal and vertical directions respectively and its $Z_{\text{c}}$ axis is aligned with the optical axis of the camera, pointing towards the scene.
Assume that the camera captures a feature point $\,^\mathcal{C} \boldsymbol{\hat{r}}$ in the image plane, which can be represented using the common pinhole camera model as:
\begin{equation}
\label{pinhole camera model}
    \,^\mathcal{C} \boldsymbol{\hat{r}} = \begin{bmatrix} \,^\mathcal{C} \hat{x} \\ \,^\mathcal{C} \hat{y} \end{bmatrix}
    = \frac{1}{\,^\mathcal{C} z} \begin{bmatrix} \,^\mathcal{C} x \\ \,^\mathcal{C} y\end{bmatrix},
\end{equation}
where the feature point $\,^\mathcal{C} \boldsymbol{\hat{r}}$ is the perspective projection of a 3-D point in the camera frame $\,^\mathcal{C} \boldsymbol{r} = [\,^\mathcal{C} x, \,^\mathcal{C} y, \,^\mathcal{C} z]^{\top}$ on the image plane.

Now consider a virtual camera frame $\mathcal{V}=\{O_{\text{v}} ~X_{\text{v}} ~Y_{\text{v}} ~Z_{\text{v}}\}$, which moves with the camera frame $\mathcal{C}$ but in such a way that the $X_{\text{v}}$--$Y_{\text{v}}$ plane is always horizontal, i.e., parallel to the $X_{\text{i}}$--$Y_{\text{i}}$ plane.
Associated with $\mathcal{V}$, let us also define a virtual image plane, whose position and orientation w.r.t. $\mathcal{V}$ is the same as the actual image plane w.r.t. $\mathcal{C}$.
The virtual image plane transformation is conducted by using the roll $\phi$ and pitch $\theta$ angles of the UAV.
The 3-D point in the camera frame $\,^\mathcal{C} \boldsymbol{r}$ is expressed in the virtual coordinate frame as $\,^\mathcal{V} \boldsymbol{r} = [\,^\mathcal{V} x, \,^\mathcal{V} y, \,^\mathcal{V} z]^{\top}$. The relationship between $\,^\mathcal{C} \boldsymbol{r}$ and $\,^\mathcal{V} \boldsymbol{r}$ can be expressed as:
\begin{equation}
\label{virtual 3D point}
    \,^\mathcal{V} \boldsymbol{r} = ^\mathcal{V}R_\mathcal{C} \,^\mathcal{C} \boldsymbol{r},
\end{equation}
where the rotation matrix $^\mathcal{V}R_\mathcal{C} = R_y(\theta) R_x(\phi)$ transforms vectors from camera frame $\mathcal{C}$ to the virtual camera frame $\mathcal{V}$.
Using \eqref{virtual 3D point} and the pinhole camera model \eqref{pinhole camera model}, virtual feature point $\,^\mathcal{V} \boldsymbol{\hat{r}}$ is computed in the virtual image plane as:
\begin{equation}
\label{virtual feature point}
    \,^\mathcal{V} \boldsymbol{\hat{r}} = \begin{bmatrix} \,^\mathcal{V} \hat{x} \\ \,^\mathcal{V} \hat{y} \end{bmatrix}
    = \frac{1}{\,^\mathcal{V} z} \begin{bmatrix} \,^\mathcal{V} x \\ \,^\mathcal{V} y\end{bmatrix} = \begin{bmatrix}
        \frac{C_\theta \,^\mathcal{C} \hat{x} + S_\theta S_\phi \,^\mathcal{C} \hat{y} + S_\theta C_\phi}{- S_\theta \,^\mathcal{C} \hat{x} + C_\theta S_\phi \,^\mathcal{C} \hat{y} + C_\theta C_\phi} \\ 
        \frac{C_\theta \,^\mathcal{C} \hat{y} - S_\phi}{- S_\theta \,^\mathcal{C} \hat{x} + C_\theta S_\phi \,^\mathcal{C} \hat{y} + C_\theta C_\phi}
    \end{bmatrix}
\end{equation}

\subsection{Problem Statement} \label{problem statement: section}
This study proposes an image-based flight controller that uses scale-independent target parameters, instead of scale-dependent state variables $\,^\mathcal{I} \boldsymbol{r}_{\text{b}}$ and $\,^\mathcal{I} \boldsymbol{v}_{\text{b}}$ given in \eqref{outer_loop_control: equation}, to maneuver the UAV towards the target without explicitly knowing the target position or heading.
These target parameters are obtained using a downward-facing monocular camera mounted onboard the UAV platform. 
Assuming the optical axis of the camera is aligned with the $Z_{\text{b}}$ axis of body frame $\mathcal{B}$, the camera continuously captures a collection of $N$ feature points of the target in the image plane.
Note that the origin of the image plane is aligned with the image center and the feature points are normalized to camera focal length $f = 1$ m.
These feature points are transformed to virtual feature points in the virtual image plane using \eqref{virtual feature point}. 
The geometric center of these virtual points is considered as the target image point in the virtual frame $\,^\mathcal{V} \hat{\boldsymbol{r}} = [\,^\mathcal{V} \hat{x}, \,^\mathcal{V} \hat{y}]^{\top}$, which is estimated as follows:
\begin{equation}
\label{image point: equation}
    \,^\mathcal{V} \hat{\boldsymbol{r}} = \begin{bmatrix} \frac{m_{10}}{m_{00}},\,\frac{m_{01}}{m_{00}} \end{bmatrix}^{\top}, 
\end{equation}
where $m_{ij}$ are the image moments given as $m_{ij} = \sum^{N}_{k=1} \,^\mathcal{V} \hat{x}_k^i \,^\mathcal{V} \hat{y}_k^j$ with the $k^{th}$ virtual feature point represented as $\,^\mathcal{V} \hat{\boldsymbol{r}}_k = [\,^\mathcal{V} \hat{x}_k, \,^\mathcal{V} \hat{y}_k]^{\top}$.

Using the common pinhole camera model \eqref{pinhole camera model}, the normalized target position $\boldsymbol{s}$ can be obtained using the target image point $\,^\mathcal{V} \hat{\boldsymbol{r}}$, given in \eqref{image point: equation}, as shown below:
\begin{equation} \label{normalized_pos: equation}
    \boldsymbol{s} = \begin{bmatrix} \,^\mathcal{V} \hat{\boldsymbol{r}} \\ 1\end{bmatrix} = \frac{\,^\mathcal{V} \boldsymbol{r}_{\text{t}}}{\,^\mathcal{V} z_{\text{t}}},
\end{equation}
where $\boldsymbol{s}$ represents the perspective projection of the target point in the virtual frame $\,^\mathcal{V} \boldsymbol{r}_{\text{t}} = [\,^\mathcal{V} x_{\text{t}}, \,^\mathcal{V} y_{\text{t}}, \,^\mathcal{V} z_{\text{t}}]^{\top}$ on the virtual image plane.
Since the onboard camera has limited capabilities, the target image point $\,^\mathcal{V} \hat{\boldsymbol{r}}$ is visible to it only when the target stays within its field of view.
Hence, $\,^\mathcal{V} \hat{\boldsymbol{r}}$ is subjected to the following visibility constraints:
\begin{equation} \label{visibility_constraints}
    -\frac{\mathcal{R}}{2f} \leq \,^\mathcal{V} \hat{\boldsymbol{r}} \leq \frac{\mathcal{R}}{2f},
\end{equation}
 where $\mathcal{R} = [r_h, r_w]^\top$ represents image resolution in pixels.

Using \eqref{target kinematics: equation} and \eqref{normalized_pos: equation}, the visual kinematics is given as:
\begin{align} \label{s_dot: equation}
    \dot{\boldsymbol{s}} &= \frac{1}{\,^\mathcal{V} z_{\text{t}}}\,^\mathcal{V} \dot{\boldsymbol{r}}_{\text{t}} - \frac{\,^\mathcal{V} \dot{z_{\text{t}}}}{\,^\mathcal{V} z_{\text{t}}^2}\,^\mathcal{V} \boldsymbol{r}_{\text{t}} = \frac{1}{\,^\mathcal{V} z_{\text{t}}} \,^\mathcal{V} \dot{\boldsymbol{r}}_{\text{t}} - \frac{(\,^\mathcal{V} \dot{\boldsymbol{r}}_{\text{t}} \cdot \hat{\boldsymbol{e}}_3)}{\,^\mathcal{V} z_{\text{t}}^2}\,^\mathcal{V} \boldsymbol{r}_{\text{t}} \nonumber \\
    &= \frac{1}{\,^\mathcal{V} z_{\text{t}}} \big( \,^\mathcal{V} \boldsymbol{v}_{\text{t/v}} - \,^\mathcal{V} \boldsymbol{\omega}_{\text{v}} \times \,^\mathcal{V} \boldsymbol{r}_{\text{t}}\big) - \bigg[\frac{1}{\,^\mathcal{V} z_{\text{t}}} \big( \,^\mathcal{V} \boldsymbol{v}_{\text{t/v}} \nonumber \\
    &~~~ - \,^\mathcal{V} \boldsymbol{\omega}_{\text{v}} \times \,^\mathcal{V} \boldsymbol{r}_{\text{t}}\big) \cdot \hat{\boldsymbol{e}}_3 \bigg] \frac{\,^\mathcal{V} \boldsymbol{r}_{\text{t}}}{\,^\mathcal{V} z_{\text{t}}} \nonumber \\
    &= \boldsymbol{h} - \,^\mathcal{V} \boldsymbol{\omega}_{\text{v}} \times \boldsymbol{s} - \left[\left( \boldsymbol{h} - \,^\mathcal{V} \boldsymbol{\omega}_{\text{v}} \times \boldsymbol{s} \right) \cdot \hat{\boldsymbol{e}}_3 \right] \boldsymbol{s},
\end{align}
where $\boldsymbol{h} = \,^\mathcal{V} \boldsymbol{v}_{\text{t/v}}/\,^\mathcal{V} z_{\text{t}}$ is the optical flow representing visual velocity of the target and  $\,^\mathcal{V} \boldsymbol{\omega}_{\text{v}} = [0, 0, -\dot{\psi}]^\top$ is the angular rate of the virtual plane in the virtual frame $\mathcal{V}$.
Given the desired normalized target position as $\boldsymbol{s}_\text{d} = [\hat{\boldsymbol{r}}_\text{d}^\top, \, 1]^\top$, the normalized position error $\boldsymbol{s}_\text{e}$ is defined as $\boldsymbol{s}_\text{e} = \boldsymbol{s} - \boldsymbol{s}_\text{d}$.
Since $\dot{\boldsymbol{s}}_\text{d} = \boldsymbol{0}_{3 \times 1}$, the visual error kinematics is derived using \eqref{s_dot: equation} as:
\begin{equation} \label{s_e_dot: equation}
    \dot{\boldsymbol{s}}_\text{e} = \dot{\boldsymbol{s}} = \boldsymbol{h} - \,^\mathcal{V} \boldsymbol{\omega}_{\text{v}} \times \boldsymbol{s} - \left[\left( \boldsymbol{h} - \,^\mathcal{V} \boldsymbol{\omega}_{\text{v}} \times \boldsymbol{s} \right) \cdot \hat{\boldsymbol{e}}_3 \right] \boldsymbol{s},
\end{equation}
Using \eqref{s_e_dot: equation}, the desired optical flow $\boldsymbol{h}_\text{d}$ is derived as:
\begin{equation} \label{h_des: equation}
    \boldsymbol{h}_\text{d} = \dot{\boldsymbol{s}}_\text{e} + \,^\mathcal{V} \boldsymbol{\omega}_{\text{v}} \times \boldsymbol{s} + \left[h_\text{rd} - \left(\,^\mathcal{V} \boldsymbol{\omega}_{\text{v}} \times \boldsymbol{s} \right) \cdot \hat{\boldsymbol{e}}_3 \right] \boldsymbol{s},
\end{equation}
where $h_\text{rd}$ is the desired radial optical flow, which is a predefined constant.
Now, in designing the outer-loop control law, the first objective is to generate suitable optical flow $\boldsymbol{h}_\text{d}$ command such that the position error $\boldsymbol{s}_\text{e}$ asymptotically converges to the origin when subjected to the visibility constraint \eqref{visibility_constraints}.

Using \eqref{target kinematics: equation}, the optical flow kinematics is given as:
\begin{align}
    \dot{\boldsymbol{h}} &= \frac{1}{\,^\mathcal{V} z_{\text{t}}}\,^\mathcal{V} \dot{\boldsymbol{v}}_{\text{t/v}} - \frac{\,^\mathcal{V} \dot{z_{\text{t}}}}{\,^\mathcal{V} z_{\text{t}}^2}\,^\mathcal{V} \boldsymbol{v}_{\text{t/v}} = \frac{1}{\,^\mathcal{V} z_{\text{t}}}\,^\mathcal{V} \dot{\boldsymbol{v}}_{\text{t/v}} - \frac{(\,^\mathcal{V} \dot{\boldsymbol{r}}_{\text{t}} \cdot \hat{\boldsymbol{e}}_3)}{\,^\mathcal{V} z_{\text{t}}^2}\,^\mathcal{V} \boldsymbol{v}_{\text{t/v}} \nonumber \\
    &= \frac{1}{\,^\mathcal{V} z_{\text{t}}} \Big[ \,^\mathcal{V} \boldsymbol{a}_{\text{t/v}} - \,^\mathcal{V} \dot{\boldsymbol{\omega}}_{\text{v}} \times \,^\mathcal{V} \boldsymbol{r}_{\text{t}} - \,^\mathcal{V} \boldsymbol{\omega}_{\text{v}} \times \left( \,^\mathcal{V} \boldsymbol{\omega}_{\text{v}} \times \,^\mathcal{V} \boldsymbol{r}_{\text{t}} \right) \nonumber \\
    &~~~ - 2 \left( \,^\mathcal{V} \boldsymbol{\omega}_{\text{v}} \times \,^\mathcal{V} \boldsymbol{v}_{\text{t/v}} \right) \Big] \nonumber \\
    &~~~ - \bigg[\frac{1}{\,^\mathcal{V} z_{\text{t}}} \big( \,^\mathcal{V} \boldsymbol{v}_{\text{t/v}} - \,^\mathcal{V} \boldsymbol{\omega}_{\text{v}} \times \,^\mathcal{V} \boldsymbol{r}_{\text{t}}\big) \cdot \hat{\boldsymbol{e}}_3 \bigg] \frac{\,^\mathcal{V} \boldsymbol{v}_{\text{t/v}}}{\,^\mathcal{V} z_{\text{t}}} \nonumber \\
    &= \frac{\,^\mathcal{V} \boldsymbol{a}_{\text{t/v}}}{\,^\mathcal{V} z_{\text{t}}} - \,^\mathcal{V} \dot{\boldsymbol{\omega}}_{\text{v}} \times \boldsymbol{s} - \,^\mathcal{V} \boldsymbol{\omega}_{\text{v}} \times \left( \,^\mathcal{V} \boldsymbol{\omega}_{\text{v}} \times \boldsymbol{s} \right) - 2 \left( \,^\mathcal{V} \boldsymbol{\omega}_{\text{v}} \times \boldsymbol{h} \right) \nonumber \\
    &~~~ - \left[\left( \boldsymbol{h} - \,^\mathcal{V} \boldsymbol{\omega}_{\text{v}} \times \boldsymbol{s} \right) \cdot \hat{\boldsymbol{e}}_3 \right] \boldsymbol{h},
\end{align}
Using $\boldsymbol{h}_\text{d}$ from \eqref{h_des: equation}, the optical flow error $\boldsymbol{h}_\text{e}$ is estimated as $\boldsymbol{h}_\text{e} = \boldsymbol{h} - \boldsymbol{h}_\text{d}$. Then, the optical flow error kinematics is given as:
\begin{align} \label{h_e_dot: equation}
    \dot{\boldsymbol{h}}_\text{e} &= \dot{\boldsymbol{h}} - \dot{\boldsymbol{h}}_\text{d} \nonumber \\
    &= \frac{\,^\mathcal{V} \boldsymbol{a}_{\text{t/v}}}{\,^\mathcal{V} z_{\text{t}}} - \,^\mathcal{V} \dot{\boldsymbol{\omega}}_{\text{v}} \times \boldsymbol{s} - \,^\mathcal{V} \boldsymbol{\omega}_{\text{v}} \times \left( \,^\mathcal{V} \boldsymbol{\omega}_{\text{v}} \times \boldsymbol{s} \right) - 2 \left( \,^\mathcal{V} \boldsymbol{\omega}_{\text{v}} \times \boldsymbol{h} \right) \nonumber \\
    &~~~ - \left[\left( \boldsymbol{h} - \,^\mathcal{V} \boldsymbol{\omega}_{\text{v}} \times \boldsymbol{s} \right) \cdot \hat{\boldsymbol{e}}_3 \right] \boldsymbol{h} - \dot{\boldsymbol{h}}_\text{d},
\end{align}

Using \eqref{quadrotor dynamics: equation} and \eqref{a_des: equation}, the relative acceleration of the target $\,^\mathcal{V}\boldsymbol{a}_{\text{t/v}}$ can be expressed in terms of the desired acceleration $^\mathcal{I}\boldsymbol{a}_{\text{d}}$ as:
\begin{align}
    \label{a_tb: equation}
    \,^\mathcal{V}\boldsymbol{a}_{\text{t/v}} &= \,^\mathcal{V}\boldsymbol{a}_{\text{t}} -\,^\mathcal{V}\boldsymbol{a}_{\text{v}} = \,^\mathcal{V}\boldsymbol{a}_{\text{t}} - \,^\mathcal{B}R_\mathcal{V}^\top \,^\mathcal{B}\boldsymbol{a}_{\text{b}} \nonumber \\
    &= \,^\mathcal{V}\boldsymbol{a}_{\text{t}} - \,^\mathcal{B}R_\mathcal{V}^\top \left(\frac{\,^\mathcal{B} \boldsymbol{F_{\text{ext}}}}{m} - \,^\mathcal{B} \boldsymbol{\omega}_{\text{b}} \times \,^\mathcal{B} \boldsymbol{v}_{\text{b}} \right) \nonumber \\
    &= - \,^\mathcal{B}R_\mathcal{V}^\top \frac{\,^\mathcal{B} \boldsymbol{F_{\text{t}}}}{m} + \,^\mathcal{V}\boldsymbol{a}_{\text{t}} - \bigg(\frac{\,^\mathcal{V} \boldsymbol{F_{\text{g}}} + \,^\mathcal{V} \boldsymbol{F_{\text{d}}}}{m} \nonumber \\
    &~~~ - \,^\mathcal{V} \boldsymbol{\omega}_{\text{b}} \times \,^\mathcal{V} \boldsymbol{v}_{\text{b}} \bigg) \nonumber \\
    &= - \,^\mathcal{I}R_\mathcal{V}^\top \, ^\mathcal{I}\boldsymbol{a}_{\text{d}} + \,^\mathcal{V}\boldsymbol{a}_{\text{t}} - \bigg(\frac{\,^\mathcal{V} \boldsymbol{F_{\text{g}}} + \,^\mathcal{V} \boldsymbol{F_{\text{d}}}}{m} \nonumber \\
    &~~~ - \,^\mathcal{V} \boldsymbol{\omega}_{\text{b}} \times \,^\mathcal{V} \boldsymbol{v}_{\text{b}} \bigg)
\end{align}

Substituting \eqref{a_tb: equation} in \eqref{h_e_dot: equation}, the perturbed error dynamics of the optical flow is given as:
\begin{equation} \label{h_e_dot_1: equation}
    \dot{\boldsymbol{h}}_\text{e} = \beta \boldsymbol{u}_\text{a} + \boldsymbol{c} + \boldsymbol{d},
\end{equation}
where 
\begin{align*}
    \boldsymbol{u}_\text{a} &= - \,^\mathcal{I}R_\mathcal{V}^\top \, ^\mathcal{I}\boldsymbol{a}_{\text{d}}, \,
    \beta = \frac{1}{z_{\text{t}}}, \\
    \boldsymbol{c} &=  - \,^\mathcal{V}\dot{\boldsymbol{\omega}}_{\text{v}} \times \boldsymbol{s} - \,^\mathcal{V}\boldsymbol{\omega}_{\text{v}} \times \left(\,^\mathcal{V} \boldsymbol{\omega}_{\text{v}} \times \boldsymbol{s} \right) - 2 \left(\,^\mathcal{V} \boldsymbol{\omega}_{\text{b}} \times \boldsymbol{h} \right) \nonumber \\
    &~~~ - \left[\left( \boldsymbol{h} - \,^\mathcal{V}\boldsymbol{\omega}_{\text{v}} \times \boldsymbol{s} \right) \cdot \hat{\boldsymbol{e}}_3 \right] \boldsymbol{h} - \dot{\boldsymbol{h}}_\text{d} \\
    \boldsymbol{d} &= \frac{\,^\mathcal{V}\boldsymbol{a}_{\text{t}}}{\,^\mathcal{V}z_{\text{t}}} - \left(\frac{\,^\mathcal{V}\boldsymbol{F_{\text{g}}} + \,^\mathcal{V}\boldsymbol{F_{\text{d}}}}{m \,^\mathcal{V}z_{\text{t}}} - \,^\mathcal{V}\boldsymbol{\omega}_{\text{b}} \times \frac{\,^\mathcal{V}\boldsymbol{v}_{\text{b}}}{\,^\mathcal{V}z_{\text{t}}} \right).
\end{align*}
Note that the function $\boldsymbol{c}$ is known and can be estimated for feedback linearization. Whereas, parameter $\beta$ and function $\boldsymbol{d}$ are unknown.

\textit{Assumption 1}: Given that $\boldsymbol{u}_\text{a}$ is bounded, the disturbance $\boldsymbol{d}$ in \eqref{h_e_dot_1: equation} is a uniformly bounded function with an unknown upper bound vector $\boldsymbol{d}_{\max}$, i.e. $\lvert \boldsymbol{d} \rvert \leq \boldsymbol{d}_{\max}$.
Moreover, the parameter $\beta$ satisfies the inequality $0 < \beta_{\min} \leq\beta\leq \beta_{\max}$, where the constants $\beta_{\min},\,\beta_{\max}$ exist but are unknown.

Now, in designing the outer-loop control law, the second objective is to generate suitable control signal $\boldsymbol{u}_\text{a}$, and hence the desired acceleration $^\mathcal{I}\boldsymbol{a}_{\text{d}}$ command, such that the optical flow error $\boldsymbol{h}_\text{e}$ asymptotically converges to the origin in the presence of uncertainty parameter $\beta$ and disturbance term $\boldsymbol{d}$.
Then, the desired acceleration $\,^\mathcal{I}\boldsymbol{a}_{\text{d}}$ is used to estimate desired roll (${\phi}_{\text{d}}$), pitch (${\theta}_{\text{d}}$), and total thrust ($\,^\mathcal{B} T_u$) commands using \eqref{euler_thrust_des: equation}.

In order to achieve orientation alignment without explicit knowledge of the target heading $\psi_t$, the proposed flight controller relies on the target image orientation $\alpha$ defined as:
\begin{equation} \label{target orientation: equation}
    \alpha = \frac{1}{2} \text{tan}^{-1}\left( \frac{2 \mu_{11}}{\mu_{20} - \mu_{02}}\right),
\end{equation}
where $\mu_{ij}$ are the centered moments given as $\mu_{ij} = \sum^{N}_{k=1} (\hat{x}_k - \hat{x})^i (\hat{y}_k - \hat{y})^j$.
Using \eqref{target orientation: equation}, the time derivative of $\alpha$ is given as:
\begin{equation}
    \dot{\alpha} = - \,^\mathcal{V} \omega_{\text{v}z}
\end{equation}
where $\,^\mathcal{V} \omega_{\text{v}z}$ is the third component of the angular rate $\,^\mathcal{V} \boldsymbol{\omega}_{\text{v}}$ of the virtual plane in the virtual frame $\mathcal{V}$. Since $\,^\mathcal{V} \boldsymbol{\omega}_{\text{v}} = [0, 0, -\dot{\psi}]^\top$, this implies that $\dot{\alpha} = \dot{\psi}$.

Hence, the outer-loop of the proposed flight controller passes the desired attitude command given by $\boldsymbol{\Phi}_{\text{d}} = [\phi_{\text{d}}, \theta_{\text{d}}, \alpha_{\text{d}}]^\top$ to the inner-loop controller to regulates the UAV attitude given by $\boldsymbol{\Phi}_{\text{b}} = [\phi_{\text{b}}, \theta_{\text{b}}, \alpha]^\top$.
Note that the inner-loop control architecture of the proposed flight controller mirrors that of the pose-based controller, given by \eqref{kinematic transformation: equation}-\eqref{tau_des: equation}, but replaces UAV heading $\psi$  with target image orientation $\alpha$  to achieve orientation alignment without explicit knowledge of the target heading $\psi_t$.

\begin{figure*}[hbt!]
\centering
    \includegraphics[width=1.0 \textwidth]{Figures/Block Diagram_v2.pdf}
    \caption{Block diagram illustrating the control pipeline for soft precise landing of a UAV on a mobile target.} 
    \label{block diagram: figure}
\end{figure*}

\section{Formulation of Control Strategy} \label{control strategy: section}
The proposed flight controller replaces the scale-dependent state variables $\,^\mathcal{I} \boldsymbol{r}_{\text{b}}$ and $\,^\mathcal{I} \boldsymbol{v}_{\text{b}}$, given in \eqref{outer_loop_control: equation}, with the scale-independent target parameters $\,^\mathcal{V} \boldsymbol{s}$ and $\,^\mathcal{V} \boldsymbol{h}$.
As mentioned in the previous section, the adoption of these target parameters poses challenges, rendering conventional PID controllers unsuitable.

\subsection{Performance Control for Visibility Constraints}
Based on the visual error kinematics \eqref{s_e_dot: equation} developed in the previous section, the prescribed performance control technique is utilized here to achieve predefined dynamic and steady-state response for visual error, as well as satisfy the visibility constraints \eqref{visibility_constraints}.
Defining image error as $\hat{\boldsymbol{r}}_\text{e} = \,^\mathcal{V} \hat{\boldsymbol{r}} - \hat{\boldsymbol{r}}_\text{d}$, the visibility constraint \eqref{visibility_constraints} can be rewritten in terms of $\hat{\boldsymbol{r}}_\text{e}$ as:
\begin{equation} \label{visibility_constraints_1}
    -\frac{\mathcal{R}}{2f} - \hat{\boldsymbol{r}}_\text{d} \leq \hat{\boldsymbol{r}}_\text{e} \leq \frac{\mathcal{R}}{2f} - \hat{\boldsymbol{r}}_\text{d},
\end{equation}
A controller achieves the prescribed performance when the image error $\hat{\boldsymbol{r}}_\text{e}$ converges to a predefined residual set that is bounded by absolutely decaying function of time. Typically, the prescribed performance can be mathematically expressed by the following inequalities:
\begin{equation}
\label{prescribed perfromance: equation}
    -\underline{C}_{k} p_{1k}(t) < \hat{r}_{\text{e}k} < \overline{C}_{k} p_{1k}(t),
\end{equation}
where $k = \{1, 2\}$, $\underline{C}_{k}$ and $\overline{C}_{k}$ are positive parameters, and the performance function $p_{1k}(t)$ is defined by:
\begin{equation} \label{performance function: equation}
    p_{1_k}(t) = \left(1 - \frac{p_{1_\infty k}}{\max{\underline{C}_{k}, \overline{C}_{k}}}\right) e^{-\xi_{1_k} t} + \frac{p_{1_\infty k}}{\max{\underline{C}_{k}, \overline{C}_{k}}},
\end{equation}
where the constant $\xi_{1_k} > 0$ denotes the convergence rate of the tracking error, and $p_{1_\infty k} > 0$ represents the steady-state performance specification, which can be set small enough to meet performance requirements.
The visibility constraint \eqref{visibility_constraints_1} is ensured by selecting $\underline{C}_{k}$ and $\overline{C}_{k}$ as:
\begin{equation}
    \underline{C} =  \frac{\mathcal{R}}{2f} + \hat{\boldsymbol{r}}_\text{d}, \, \overline{C} =  \frac{\mathcal{R}}{2f} - \hat{\boldsymbol{r}}_\text{d}, 
\end{equation}

To achieve target tracking with prescribed performance \eqref{prescribed perfromance: equation}, an error transformation first is introduced. The transformation function $\mathcal{S}_1(\boldsymbol{\zeta}_1)$ is defined such that
\begin{equation}
    \label{error transformation: equation}
    \hat{\boldsymbol{r}}_\text{e} = \mathcal{S}_1(\boldsymbol{\zeta}_1) \boldsymbol{p}_1(t), \, \mathcal{S}_1(\boldsymbol{\zeta}_1) = \text{diag}(S_1(\zeta_{1_k})),
\end{equation}
where function $S_1(\zeta_{1_k})$ is strictly increasing and satisfies the following requirements:
\begin{align*}
&-\underline{C}_{k} < S_1(\zeta_{1_k}) < \overline{C}_{k}, \\ 
&\lim_{\zeta_{1_k} \to \infty} S_1(\zeta_{1_k}) = \overline{C}_{k}, \\
&\lim_{\zeta_{1_k} \to -\infty} S_1(\zeta_{1_k}) = -\underline{C}_{k}.
\end{align*}
According to the above mentioned requirements, the error transformation function $S_1(\zeta_{1_k})$ can be chosen as:
\begin{equation}
    \label{transformation function: equation}
    S_1(\zeta_{1_k}) = \left(\frac{e^{\zeta_{1_k}} \overline{C}_{k} - \underline{C}_{k}}{e^{\zeta_{1_k}} + 1} \right).
\end{equation}
Since $S_1(\zeta_{1_k})$ is strictly increasing, its inverse function always exists. 
The inverse function $S_1^{-1}(\hat{r}_{\text{e}k}, p_{1_k})$ can be described by:
\begin{equation} \label{inverse transformation function: equation}
    S_1^{-1}(\hat{r}_{\text{e}k}, p_{1_k}) = \zeta_{1_k} = \ln{\left(\frac{\hat{r}_{\text{e}k} + \underline{C}_{k} p_{1_k}(t)}{\overline{C}_{k} p_{1k}(t) - \hat{r}_{\text{e}k}} \right)}.
\end{equation}
The time-derivative of transformed error $\zeta_{1_k}$ is given as:
\begin{equation} \label{transformed error derivative: equation}
    \dot{\boldsymbol{\zeta}}_1 = \mathcal{G}_1 \big[ \dot{\hat{\boldsymbol{r}}}_{\text{e}} - \mathcal{S}_1(\boldsymbol{\zeta}_1) \dot{\boldsymbol{p}}_1(t)\big],    
\end{equation}
where
\begin{align*}
    \mathcal{G}_1 &= \left[\boldsymbol{p}_1(t)^{\top} \frac{\partial \mathcal{S}_1(\boldsymbol{\zeta}_1)}{\partial \boldsymbol{\zeta}_1}\right]^{-1} = \text{diag}(g_{1_k}(\zeta_{1_k}, p_{1_k})), \nonumber \\
    g_{1_k}(\zeta_{1_k}, p_{1_k}) &= \frac{(e^{\zeta_{1_k}}+1)^2}{2\,e^{\zeta_{1_k}} p_{1_k}(t)}.
\end{align*}

Now, the the visual error kinematics on the unconstrained error is considered. Hence, using \eqref{transformed error derivative: equation}, $\dot{\boldsymbol{s}}_\text{e}$ can be analytically computed as:
\begin{equation} \label{s_e_d_dot: equation}
    \dot{\boldsymbol{s}}_\text{e} = \begin{bmatrix} \dot{\hat{\boldsymbol{r}}}_{\text{e}} \\ 0 \end{bmatrix} = \begin{bmatrix} \mathcal{G}_1^{-1} \dot{\boldsymbol{\zeta}}_{1} + \mathcal{S}_1(\boldsymbol{\zeta}_1) \dot{\boldsymbol{p}}_1(t) \\ 0 \end{bmatrix}.   
\end{equation}
Substituting \eqref{s_e_d_dot: equation} in \eqref{h_des: equation}, we obtain
\begin{align} \label{h_des_1: equation}
    \boldsymbol{h}_\text{d} &= \begin{bmatrix} \mathcal{G}_1^{-1} \dot{\boldsymbol{\zeta}}_{1} + \mathcal{S}_1(\boldsymbol{\zeta}_1) \dot{\boldsymbol{p}}_1(t) \\ 0 \end{bmatrix} - \,^\mathcal{V} \boldsymbol{\omega}_{\text{v}} \times \boldsymbol{s} \nonumber \\
    &~~~ - \left[h_\text{rd} - \left(\,^\mathcal{V} \boldsymbol{\omega}_{\text{v}} \times \boldsymbol{s} \right) \cdot \hat{\boldsymbol{e}}_3 \right] \boldsymbol{s},
\end{align}

So, the first objective translates to the formulation of a control law that synthesizes suitable optical flow $\boldsymbol{h}_\text{d}$ command such that the unconstrained error $\boldsymbol{\zeta}_{1}$ asymptotically converges to the origin. To this end, a simple PID control law is proposed as given below:
\begin{equation} \label{zeta_e_dot_d: equation}
    \dot{\boldsymbol{\zeta}}_{1} = - K_{\zeta p} \boldsymbol{\zeta}_{1} - K_{\zeta i} \int \boldsymbol{\zeta}_{1} dt - K_{\zeta d} \frac{d\boldsymbol{\zeta}_{1}}{dt}.
\end{equation}
Substituting \eqref{zeta_e_dot_d: equation} in \eqref{h_des_1: equation}, we obtain:
\begin{align} \label{h_des_2: equation}
    \boldsymbol{h}_\text{d} &= \begin{bmatrix} - \mathcal{G}_1^{-1} \left(K_{\zeta p} \boldsymbol{\zeta}_{1} + K_{\zeta i} \int \boldsymbol{\zeta}_{1} dt + K_{\zeta d} \frac{d\boldsymbol{\zeta}_{1}}{dt}\right) + \mathcal{S}_1(\boldsymbol{\zeta}_1) \dot{\boldsymbol{p}}_1(t) \\ 0 \end{bmatrix} \nonumber \\
    &~~~ - \,^\mathcal{V} \boldsymbol{\omega}_{\text{v}} \times \boldsymbol{s} - \left[h_\text{rd} - \left(\,^\mathcal{V} \boldsymbol{\omega}_{\text{v}} \times \boldsymbol{s} \right) \cdot \hat{\boldsymbol{e}}_3 \right] \boldsymbol{s}.
\end{align}

\subsection{Performance Function for Optical Flow}
Based on the perturbed error dynamics \eqref{h_e_dot_1: equation} developed in the previous section, the prescribed performance control technique is utilized here to achieve predefined dynamic and steady-state response for optical flow error, as well as ensure the boundedness during landing.
To this end, the optical flow error $\boldsymbol{h}_{e}$ is constrained as:
\begin{equation}
    h_{e_k} \in (- p_{2_k}(t), \, p_{2_k}(t)),
\end{equation}
where $k = \{1, 2, 3\}$ and the performance function $\boldsymbol{p}_{2}(t))$ is defined by:
\begin{equation}
    \boldsymbol{p}_2(t) = e^{-\Xi_i t} (\boldsymbol{p}_{2_0} - \boldsymbol{p}_{2_\infty}) + \boldsymbol{p}_{2_\infty}   
\end{equation}
such that
$\lim_{\text{t} \to \infty} p_{2_k}(t) = p_{2_\infty k} > 0$. Here, $p_{2_0 k}$, $p_{2_\infty k}$, and matrix $\Xi_2 = \text{diag}(\xi_{2_k})$ are predefined positive constants, which represent the maximum allowable overshoot, the steady-state residual set and the minimum convergence rate of $\boldsymbol{h}_{e}$, respectively.

A non-linear transformation is formulated to map the constrained output error $\boldsymbol{h}_{e}$ onto a new unconstrained variable $\boldsymbol{\zeta}_{2}$ such that:
\begin{align}
    \label{optical flow error transformation: equation}
    \boldsymbol{h}_{e} &= \mathcal{S}_2(\boldsymbol{\zeta}_2) \boldsymbol{p}_2(t), \, \mathcal{S}_2(\boldsymbol{\zeta}_2) = \text{diag}(S_2(\zeta_{2_k})), \nonumber \\
    S_2(\zeta_{2_k}) &= \left(\frac{e^{\zeta_{2_k}} - 1}{e^{\zeta_{2_k}} + 1} \right)
\end{align}
with $S_2(\zeta_{2_k})$ possessing the following properties:
\begin{enumerate}
    \item[I.] $ -1 < S_2(\zeta_{2_k}) < 1 $, $\forall~\zeta_{2_k}(t) \in \mathbb{R},$ 
    \item[II.] $\lim_{\zeta_{2_k} \to \infty} S_2(\zeta_{2_k}) = 1$,
    \item[III.] $\lim_{\zeta_{2_k} \to -\infty} S_2(\zeta_{2_k}) = -1$
\end{enumerate}

Moreover, the components of the unconstrained output $\boldsymbol{\zeta}_2$ can be expressed as:
\begin{equation} \label{inverse transformation function for inner loop}
    \zeta_{2_k} = S^{-1}_2(h_{e_k}, p_{2_k}) = ln \left(\frac{p_{2_k} + h_{e_k}}{p_{2_k} - h_{e_k}} \right).
\end{equation}
Note that the inverse transformation $\mathcal{S}_2^{-1}: (-1, 1) \to \mathbb{R}$ is well-defined, provided $\zeta_{2_k}(t) \in \mathbb{R}$.

Differentiating Eq. \eqref{optical flow error transformation: equation} and rearranging, we have:
\begin{equation} \label{transformed error derivative: equation 2}
    \dot{\boldsymbol{\zeta}}_2 = \mathcal{G}_2 \big[\dot{\boldsymbol{h}}_{e} - \mathcal{S}_2(\boldsymbol{\zeta}_2) \dot{\boldsymbol{p}}_2(t)\big],    
\end{equation}
where
\begin{align*}
    \mathcal{G}_2 &= \left[\boldsymbol{p}_2(t)^{\top} \frac{\partial \mathcal{S}_2(\boldsymbol{\zeta}_2)}{\partial \boldsymbol{\zeta}_2}\right]^{-1} = \text{diag}(g_{2_k}(\zeta_{2_k}, p_{2_k}), \nonumber \\
    g_{2_k}(\zeta_{2_k}, p_{2_k}) &= \frac{(e^{\zeta_{2_k}}+1)^2}{2\,e^{\zeta_{2_k}} p_{2_k}(t)}
\end{align*}
Substituting \eqref{h_e_dot_1: equation} in the above equation \eqref{transformed error derivative: equation 2} to derive an unconstrained perturbed system equation enveloped with prescribed performance as shown below:
\begin{equation} \label{unconstrained system: equation}
    \dot{\boldsymbol{\zeta}}_2 = \mathcal{G}_2 \big[\beta \boldsymbol{u}_\text{a} + \boldsymbol{c} + \boldsymbol{d} - \mathcal{S}_2(\boldsymbol{\zeta}_2) \dot{\boldsymbol{p}}_2(t) \big]
\end{equation}

\subsection{Adaptive Control Law} \label{adaptive controller for inner loop : section}
The second objective translates to the formulation of a control law that synthesizes suitable control signal $\boldsymbol{u}_\text{a}(t)$ to drive the non-linear map-based transformed output $\boldsymbol{\zeta}_2$ to a small uniform ultimate bound (UUB) around the origin in finite time through online estimation and rejection of unknown bounded uncertainty parameters $\beta$, $\boldsymbol{d}$ shown in Eq. \eqref{unconstrained system: equation}.
To this end, an adaptive first-order sliding mode control (AFOSMC) strategy is designed that compensates for the system uncertainty without any a priori knowledge of its upper bound.
This control strategy accounts for non-linearity and singularity inherent to optical flow-guided landing and alleviates the adverse effects of actuator faults and input saturation.

Similar to the outer loop controller, the following first-order PI sliding surface $\boldsymbol{\sigma}$ is defined for the inner loop output error vector $\boldsymbol{\zeta}_2$:
\begin{equation} \label{sigma: equation}
    \boldsymbol{\sigma} = \boldsymbol{\zeta}_2 + \mathcal{X} \int_{\text{t}_0}^t \boldsymbol{\zeta}_2~d\tau,
\end{equation}
where the sliding surface $\sigma_{k}$ represents the exponentially converging dynamics of output error $\zeta_{2_k}(t)$ with a decaying rate $\chi_{2_k} > 0$ for $k = \{1, 2, 3\}$ such that $\mathcal{X} = \text{diag}(\chi_{k})$.
The time-derivative of $\boldsymbol{\sigma}$ can be computed as:
\begin{equation} \label{sigma derivative: equation 1}
    \dot{\boldsymbol{\sigma}} = \dot{\boldsymbol{\zeta}}_2 + \mathcal{X}~\boldsymbol{\zeta}_2
\end{equation}
The following equation can be obtained by substituting  \eqref{unconstrained system: equation} in the above Eq. \eqref{sigma derivative: equation 1}:
\begin{align}
    \label{sigma derivative: equation 2}
    \dot{\boldsymbol{\sigma}} =& ~\mathcal{G}_2 \big[\beta \boldsymbol{u}_\text{a} + \boldsymbol{c} + \boldsymbol{d} - \mathcal{S}_2(\boldsymbol{\zeta}_2) \dot{\boldsymbol{p}}_2(t) \big] + \mathcal{X}~\boldsymbol{\zeta}_2 \nonumber \\
    = &~ \beta \mathcal{G}_2 \big[\boldsymbol{u}_\text{a} +  \boldsymbol{c} - \mathcal{S}_2(\boldsymbol{\zeta}_2) \dot{\boldsymbol{p}}_2(t) + \mathcal{G}_2^{-1} \mathcal{X}\,\boldsymbol{\zeta}_2 \big] + \beta_{\min} \, \mathcal{D} \boldsymbol{\theta}, 
\end{align}
where 
\begin{align*}
    \mathcal{D} &= \beta_{\min}^{-1} \, \begin{bmatrix} \beta - 1 & \lVert \mathcal{G}_2 \rVert^{-1} \mathcal{G}_2 \boldsymbol{d} \end{bmatrix}, \\
    \boldsymbol{\theta} &= \begin{bmatrix}
    \mathcal{G}_2 \big[- \boldsymbol{c} + \mathcal{S}_2(\boldsymbol{\zeta}_2) \dot{\boldsymbol{p}}_2(t) - \mathcal{G}_{2}^{-1} \mathcal{X}\,\boldsymbol{\zeta}_2 \big] \\ \lVert \mathcal{G}_2 \rVert
    \end{bmatrix},
\end{align*}
There exists an unknown positive constant $\tilde{d} = \beta_{\min}^{-1} \Big[\text{max}\{|\beta_{\max}{-}1|,|\beta_{\min}{-}1|\} + \lVert \boldsymbol{d}_{\max} \rVert \Big]$ such that $\lVert\mathcal{D}\rVert\leq\tilde{d}$.
 
Finally, the following control law has been proposed:
\begin{align} \label{adaptive control law: equation}
    \boldsymbol{u}_\text{a} =&~ \mathcal{G}_2^{-1} \big[ \boldsymbol{u}_{\text{sw}} + \boldsymbol{u}_{\text{eq}} \big], \nonumber \\
    \boldsymbol{u}_{\text{sw}} =& - \Gamma \boldsymbol{\sigma} - \lVert \boldsymbol{\theta} \rVert \, \text{sat}(\mathcal{E}^{-1}\boldsymbol{\sigma}) \boldsymbol{\kappa}(t)\nonumber\\
    \boldsymbol{u}_{\text{eq}} =&~\mathcal{G}_2 \big[- \boldsymbol{c} + \mathcal{S}_2(\boldsymbol{\zeta}_2) \dot{\boldsymbol{p}}_2(t) - \mathcal{G}_{2}^{-1} \mathcal{X}\,\boldsymbol{\zeta}_2 \big],
\end{align}
where $\Gamma = \text{diag}(\gamma_{k})$  is the proportional control matrix with positive constants $\gamma_{k}$ and $\boldsymbol{\kappa}(t)$ is the adaptive sliding mode control parameter.
The design matrix $\mathcal{E}_i$ is defined as $\mathcal{E} = \text{diag}(\varepsilon_{k})$ such that $0 < \varepsilon_{k} << 1$, where the design constants $\varepsilon_{k}$ represent the sensitivity of the smooth transition across the sliding surface for sliding mode control.
the sliding surface for sliding mode control.
The saturation matrix $\text{sat}(\mathcal{E}^{-1}\boldsymbol{\sigma}) = \text{diag}(\text{sat}(\sigma_{k}/\varepsilon_{k}))$, where the saturation function sat($\sigma_k/\epsilon_k$) is defined as
\begin{align}
\label{sat_fn}
\text{sat}(\sigma_k/\epsilon_k)=\begin{cases}
    \sigma_k/\epsilon_k,~~\,\text{if}~~\rvert\sigma_k\rvert\,<\,\epsilon_k\\
    \lceil \sigma_k\rfloor^0,~\text{otherwise}
    \end{cases}
\end{align}
with $\lceil \sigma_k \rfloor^0$ representing the sign function.
The adaptive gain $\boldsymbol{\kappa}(t)$ is updated as:
\begin{equation} \label{adaptive law: equation}
    \dot{\boldsymbol{\kappa}}(t) =  \lVert \boldsymbol{\theta} \rVert \mathcal{N} \lvert \boldsymbol{\sigma} \rvert - \mathcal{N} \mathcal{P} \boldsymbol{\kappa}(t), \, \boldsymbol{\kappa}(0) \in \mathbb{R}_+^3,
\end{equation}
where $\lvert \boldsymbol{\sigma} \rvert  = \left[\lvert \sigma_{1} \rvert,\, \lvert \sigma_{2} \rvert,\, \lvert \sigma_{3} \rvert\right]^{\top}$. Also, $\mathcal{N} = \text{diag}(\eta_{k})$ and $\mathcal{P} = \text{diag}(\rho_{k})$ are positive gain matrices.

Eq. \eqref{adaptive control law: equation} is incorporated in Eq. \eqref{sigma derivative: equation 2} to obtain the following expression:
\begin{align}
    \label{sigma derivative: equation 3}
    \dot{\boldsymbol{\sigma}} =&~ - \beta \big[\Gamma \boldsymbol{\sigma} + \lVert \boldsymbol{\theta} \rVert \, \text{sat}( \mathcal{E}^{-1}\boldsymbol{\sigma}) \boldsymbol{\kappa}(t)\big] + \beta_{\min} \, \mathcal{D} \boldsymbol{\theta},
\end{align}

\subsection{Lyapunov Stability Analysis}
The stability analysis starts with the following definition of global uniform ultimate boundedness of the signal $\boldsymbol{\chi}(t)$:

\textit{Definition 1}:
 Consider a multi-variable signal $\boldsymbol{\chi}(t) \in \mathbb{R}^n, ~n \in [0, \infty)$. Assume that $\exists$ $\lambda, ~ \lambda_{\text{u}} \in \mathbb{R}_+$, where $\lambda, ~ \lambda_{\text{u}}$ are positive constants and $\forall \, 0 < \lambda < \lambda_{\text{u}}$, $\exists ~ \mathcal{T} \geq 0$ and $\Lambda > 0$ such that
\begin{equation} \label{GUUB: Definition}
    \lVert \boldsymbol{\chi}(0) \rVert \leq \lambda ~\implies~ \lVert \boldsymbol{\chi}(t) \rVert \leq \Lambda,  \, \forall \, t \geq \mathcal{T}
\end{equation}
\noindent and $\lambda$ can be arbitrarily large. Then, the signal $\boldsymbol{\chi}(t)$ is said to be \textit{global uniformly ultimately bounded} by the ultimate bound $\Lambda$.

\textit{Lemma 1}:
Let $V$ be a continuously differentiable function such that 
\begin{align}
    k_1 \lVert \boldsymbol{\chi} \rVert^2_2 \leq V(\boldsymbol{\chi}) \leq k_2 \lVert \boldsymbol{\chi} \rVert^2_2, ~ \dot{V}(\boldsymbol{\chi}) \leq - k_3 V(\boldsymbol{\chi}),\\ \forall \, \lVert \boldsymbol{\chi} \rVert \geq \varkappa \nonumber
\end{align}
\noindent with $k_1$, $k_2$, $k_3 > 0$. Consider a ball $\mathcal{B}_r$ of radius $r$ and let $\varkappa \leq r (k_1/k_2)^{\frac{1}{2}} \triangleq \lambda$.
\noindent Then, for every initial state $\boldsymbol{\chi}(0)$ satisfying $\lVert \boldsymbol{\chi}(0) \rVert^2_2 \leq \lambda$, $\exists$ $\mathcal{T} \geq 0$, such that \eqref{GUUB: Definition} holds true with $\Lambda \triangleq \varkappa (k_2/k_1)^{\frac{1}{2}}$ and an arbitrarily large $\lambda$. 

\textit{Theorem 1}:
Consider the perturbed nonlinear system \eqref{h_e_dot_1: equation} such that the initial value of the optical flow error $\boldsymbol{h}_{e}(0)$ is within the constrained region, i.e., $- \boldsymbol{p}_{2_0} < \boldsymbol{h}_{e}(0) < \boldsymbol{p}_{2_0}$. Then, under \textit{Assumption 1}, system \eqref{unconstrained system: equation} under the proposed control law \eqref{adaptive control law: equation} accomplishes exponential convergence of $\boldsymbol{\zeta}_2$ to a UUB given by
\begin{equation}
    \vartheta = \sqrt{\frac{3 \beta_{\min}\lambda_{\max}(\mathcal{P})}{\left(\varphi_1 - \varphi_2\right)}} \tilde{d},
\end{equation}
where $\varphi_1 \triangleq \beta_{\min} \frac{\text{ min}\big(2 \lambda_{\min}\left(\Gamma \right),\lambda_{\min}\left(\mathcal{P} \right)\big)}{\text{max}\big(1,\beta_{\min} \lambda_{\max}\left(\mathcal{N}^{-1} \right)\big)}>0$, $0<\varphi_2<\varphi_1$.
Moreover, the control law \eqref{adaptive control law: equation} ensures that the trajectory of the system \eqref{h_e_dot_1: equation} lies within the prescribed constraint sets $- \boldsymbol{p}_{2}(t) < \boldsymbol{h}_{e} < \boldsymbol{p}_{2}(t) \, \forall \, t \geq 0$.

\textit{Proof:}
As a first step, observe that $\zeta_{2_k}(0) = S_{2_k}^{-1}(h_{e_k}(0) / p_{2_0 k})$ is well-defined for $\lvert h_{e_k}(0) \rvert < p_{2_0 k}$ for $k = \{1, 2, 3\}$. The dynamics of $\boldsymbol{\zeta}_2$ then evolves according to the system \eqref{unconstrained system: equation} under the control law \eqref{adaptive control law: equation} of the PLASMC strategy.

\noindent The following candidate Lyapunov function $V(t)$ is now considered for subsequent analysis:
\begin{align} \label{Lyapunov function: equation}
    V(t) =&~ \frac{1}{2}\boldsymbol{\sigma}^{\top}\boldsymbol{\sigma} + \frac{\beta_{\min}}{2} \left[\boldsymbol{\kappa}(t) - \tilde{d} \boldsymbol{1}\right]^{\top}  \mathcal{N}^{-1} \left[\boldsymbol{\kappa}(t) - \tilde{d} \boldsymbol{1}\right],
\end{align}
where the unknown positive constant $\tilde{d}$ satisfies $\lVert\mathcal{D}\rVert \leq \tilde{d}$ and $\boldsymbol{1} = [1,\,1,\,1]^{\top}$.

Using \eqref{sigma derivative: equation 3} and \eqref{adaptive law: equation}, the time-derivative of \eqref{Lyapunov function: equation} can be written as
\begin{align}
\label{Lyapunov derivative: equation 1}
    \dot{V}(t)=~&\boldsymbol{\sigma}^{\top} \dot{\boldsymbol{\sigma}} + \beta_{\min} \left[\boldsymbol{\kappa}(t) - \tilde{d} \boldsymbol{1}\right]^{\top} \mathcal{N}^{-1} \dot{\boldsymbol{\kappa}}(t)\nonumber\\
    = & - \beta \, \boldsymbol{\sigma}^{\top} \Big[\Gamma \boldsymbol{\sigma} + \lVert \boldsymbol{\theta} \rVert \, \text{sat}( \mathcal{E}^{-1}\boldsymbol{\sigma}) \boldsymbol{\kappa}(t)\Big] + \beta_{\min} \, \boldsymbol{\sigma}^{\top} \mathcal{D} \boldsymbol{\theta} \nonumber\\
    & + \beta_{\min} \left[\boldsymbol{\kappa}(t) - \tilde{d} \boldsymbol{1}\right]^{\top} \mathcal{N}^{-1} \Big[\lVert \boldsymbol{\theta} \rVert \mathcal{N} \lvert \boldsymbol{\sigma} \rvert - \mathcal{N} \mathcal{P} \boldsymbol{\kappa}(t) \Big].
\end{align} 
Since the diagonal matrix $\Gamma$ is positive-definite, it can be shown that $\boldsymbol{\sigma}^{\top} \Gamma \boldsymbol{\sigma} > 0$.
Moreover, $\boldsymbol{\sigma}^{\top} \text{sat}( \mathcal{E}^{-1}\boldsymbol{\sigma}) \approx \lvert \boldsymbol{\sigma} \rvert^{\top} \in \mathbb{R}_+^3$ and $\boldsymbol{\kappa}(t) \in \mathbb{R}_+^3$ from Eq. \eqref{adaptive law: equation}.
% Therefore, $\boldsymbol{\sigma}^{\top} \text{sat}( \mathcal{E}^{-1}\boldsymbol{\sigma}) \boldsymbol{\kappa}(t) > 0$.
Since $\beta \geq \beta_{\min} > 0$, \eqref{Lyapunov derivative: equation 1} can be rewritten as:
\begin{align}
\label{Lyapunov derivative: equation 2}
    \dot{V}(t)~{\leq}~& - \beta_{\min} \boldsymbol{\sigma}^{\top} \Gamma \boldsymbol{\sigma} - \beta_{\min} \lVert \boldsymbol{\theta} \rVert \lvert \boldsymbol{\sigma} \rvert^{\top} \boldsymbol{\kappa}(t) \nonumber \\
    & + \beta_{\min} \lVert\mathcal{D}\rVert \lVert \boldsymbol{\theta} \rVert \lvert \boldsymbol{\sigma} \rvert^{\top} \boldsymbol{1} \nonumber \\
    & + \beta_{\min} \left[\boldsymbol{\kappa}(t) - \tilde{d} \boldsymbol{1}\right]^{\top} \left[\lVert \boldsymbol{\theta} \rVert \lvert \boldsymbol{\sigma} \rvert - \mathcal{P} \boldsymbol{\kappa}(t) \right] \nonumber \\
    \leq ~ & - \beta_{\min} \boldsymbol{\sigma}^{\top} \Gamma \boldsymbol{\sigma} - \beta_{\min} \left( \tilde{d} - \lVert\mathcal{D}\rVert \right) \lVert \boldsymbol{\theta} \rVert \lvert \boldsymbol{\sigma} \rvert^{\top} \boldsymbol{1} \nonumber \\
    & - \beta_{\min} \left[\boldsymbol{\kappa}^{\top}(t) \mathcal{P} \boldsymbol{\kappa}(t) - \tilde{d} \boldsymbol{1}^{\top} \mathcal{P} \boldsymbol{\kappa}(t) \right]
\end{align}
Given that $\lVert\mathcal{D}\rVert \leq \tilde{d}$, it can be shown that $\beta_{\min} \left( \tilde{d} - \lVert\mathcal{D}\rVert \right) \lVert \boldsymbol{\theta} \rVert \lvert \boldsymbol{\sigma} \rvert^{\top} \boldsymbol{1} \geq 0$.
Moreover, since $\mathcal{P}$ is a positive-definite matrix, $\beta_{\min} \left[\boldsymbol{\kappa}^{\top}(t) \mathcal{P} \boldsymbol{\kappa}(t) - \tilde{d} \boldsymbol{1}^{\top} \mathcal{P} \boldsymbol{\kappa}(t) \right] > \frac{\beta_{\min}}{2} \left\{\left[\boldsymbol{\kappa}(t) - \tilde{d} \boldsymbol{1}\right]^{\top} \mathcal{P} \left[\boldsymbol{\kappa}(t) - \tilde{d} \boldsymbol{1}\right] - \tilde{d}^2 \boldsymbol{1}^{\top} \mathcal{P} \boldsymbol{1} \right\}$.
Hence,
\begin{align}
\label{Lyapunov derivative: equation 3}
\dot{V}(t) ~ \leq ~ &-\beta_{\min}\lambda_{\min}(\Gamma) \lVert \boldsymbol{\sigma} \rVert^2 \nonumber\\
& - \frac{\beta_{\min} }{2}\left[\boldsymbol{\kappa}(t) - \tilde{d} \boldsymbol{1}\right]^{\top} \mathcal{P} \left[\boldsymbol{\kappa}(t) - \tilde{d} \boldsymbol{1}\right] \nonumber\\
& +\frac{\beta_{\min} }{2}\tilde{d}^2 \boldsymbol{1}^{\top} \mathcal{P} \boldsymbol{1} \nonumber \\
\leq ~ &-\beta_{\min} \Big[ \lambda_{\min}(\Gamma) \lVert \boldsymbol{\sigma} \rVert^2 + \frac{1}{2}\lambda_{\min}(\mathcal{P}) \lVert \boldsymbol{\kappa}(t) - \tilde{d} \boldsymbol{1} \rVert^2 \nonumber\\
& - \frac{3}{2} \lambda_{\min}(\mathcal{P}) \tilde{d}^2 \Big] \nonumber \\
\leq~ &-\varphi_1 V(t) + \frac{3}{2} \beta_{\min} \lambda_{\max}(\mathcal{P}) \tilde{d}^2,
\end{align}
where $\varphi_1 \triangleq \beta_{\min} \frac{\text{ min}\big(2 \lambda_{\min}\left(\Gamma \right),\lambda_{\min}\left(\mathcal{P} \right)\big)}{\text{max}\big(1,\beta_{\min} \lambda_{\max}\left(\mathcal{N}^{-1} \right)\big)}>0$.
For the choice of a scalar $\varphi_2$ satisfying $0<\varphi_2<\varphi_1$, the above inequality \eqref{Lyapunov derivative: equation 3} is rewritten as:
\begin{align}
\label{Lyapunov derivative: equation 4}
\dot{V}(t)\leq~&-\varphi_2 V(t)-(\varphi_1-\varphi_2) V(t) + \frac{3}{2} \beta_{\min} \lambda_{\max}(\mathcal{P}) \tilde{d}^2 \nonumber\\
\leq ~&-\varphi_2 V(t), \, \forall\, V(t)\geq \Theta=\frac{3 \beta_{\min} \lambda_{\max}(\mathcal{P}) \tilde{d}^2}{2(\varphi_1-\varphi_2)}
\end{align}

Invoking \textit{Lemma 1}, any trajectory enters the ball $\mathcal{B}_{r} = \left \{ \left ( \boldsymbol{\sigma}, \, \boldsymbol{\kappa} \right ) \in \mathbb{R}^8 | V \leq \Theta \right\}$ exponentially in finite-time and stays in $\mathcal{B}_{r}$ for all future time. Note that the size of the ball can be made suitably small through a suitable selection of the control parameters $\lambda_{\min}(\mathcal{P})$, $\lambda_{\max}(\mathcal{P})$, $\lambda_{\max}(\mathcal{N}^{-1}) $ and $\lambda_{\min}(\Gamma)$, so that by recognizing that $V(t) \geq \lVert \boldsymbol{\sigma} \rVert^2/2$, one has the ultimate bound for $\boldsymbol{\sigma}$ as $\lVert \boldsymbol{\sigma} \rVert \leq \vartheta = \left[3\beta_{\min}\lambda_{\max}(\mathcal{P})/\left(\varphi_1 - \varphi_2\right)\right]^{\frac{1}{2}}  \tilde{d}$. This bound is uniform as it is independent of initial conditions. 
Consequently, it follows from \eqref{sigma: equation} that the output error variable $\boldsymbol{\zeta}_2$ converges to a small bound around the origin so that from \eqref{optical flow error transformation: equation}, $- \boldsymbol{p}_{2}(t) < \boldsymbol{h}_{e} < \boldsymbol{p}_{2}(t)$ is satisfied. Thus the prescribed performance constraint is ensured $\forall \, t \geq 0$. This ends the proof.

\section{Numerical Simulation Results}
This section provides numerical simulation results for soft precise landing on both stationary and mobile platform.

The control parameters are tabulated in \textit{Table} \ref{control params: table}. For all the test cases, the initial position of the UAV is $\,^\mathcal{I} \boldsymbol{r}_{\text{b}}(0) = [1.0, -1.0, 5.0]^\top$ m and the initial heading is $\psi(0) = 10.0^\circ$.
Each figure has been configured as given below:
\begin{enumerate}
    \item The first two rows show state, input and output plots regarding the UAV lateral control.
    \item The third row show plots regarding the UAV vertical control.
    \item The last two row shows plots regarding the UAV lateral control.
\end{enumerate}

\begin{table}
    \centering
    \begin{tabular}{|c|c|}
    \hline
        \textbf{Control} & \multirow{2}{*}{\textbf{Control Values}}\\
         \textbf{Parameter} & \\ \hline
         \multicolumn{2}{|c|}{\textbf{Performance Function for Visibility Constraint}} \\ \hline
         $\Xi_1$ & $\text{diag}([0.1, 0.1])$\\
         $\boldsymbol{p}_{1_\infty}$ & $[0.1, 0.1]^\top$ \\ [0.3em] \hline 
         \multicolumn{2}{|c|}{\textbf{PID-based Normalized Position Control}} \\ \hline
         $K_{\zeta p}$ & $\text{diag}([2.0, 2.0])$ \\
         $K_{\zeta i}$ & $\text{diag}([2.0, 2.0])$ \\
         $K_{\zeta d}$ & $\text{diag}([0.1, 0.01])$ \\ [0.3em] \hline 
         \multicolumn{2}{|c|}{\textbf{Performance Function for Optical Flow}} \\ \hline
         $\Xi_2$ & $\text{diag}([0.2, 0.2, 0.3])$ \\
         $\boldsymbol{p}_{2_0}$ & $[50.0, 50.0, 10.0]^\top$ \\
         $\boldsymbol{p}_{2_\infty}$ & $[0.1, 0.1, 0.5]^\top$ \\ [0.3em] \hline
         \multicolumn{2}{|c|}{\textbf{ASMC-based Optical Flow Control}} \\ \hline
         $\Omega$ & $\text{diag}([0.01, 0.12, 0.1])$ \\
         $\Gamma$ & $\text{diag}([0.0003, 0.0008, 4.2])$ \\
         $P$ & $\text{diag}([0.02, 0.01, 0.01])$ \\ 
         $N$ & $\text{diag}([0.01, 0.01, 0.02])$ \\ 
         $\boldsymbol{\kappa}$ & $\text{diag}([0.003, 0.01, 0.002])$ \\
         $E$ & $\text{diag}([0.8, 0.5, 0.6])$ \\ [0.3em] \hline
         \multicolumn{2}{|c|}{\textbf{PID-based Attitude Control}} \\ \hline
         $K_{\Phi p}$ & $\text{diag}([1.0, 9.0, 5.0])$ \\
         $K_{\Phi i}$ & $\text{diag}([0.004, 0.01, 0.01])$ \\
         $K_{\Phi d}$ & $\text{diag}([0.03, 0.3, 0.2])$ \\ [0.3em] \hline
         \multicolumn{2}{|c|}{\textbf{PID-based Attitude Rate Control}} \\ \hline
         $K_{\omega p}$ & $\text{diag}([1.0, 9.0, 2.0])$ \\
         $K_{\omega i}$ & $\text{diag}([0.003, 0.1, 1.4])$ \\
         $K_{\omega d}$ & $\text{diag}([0.1, 0.1, 0.1])$ \\ \hline
    \end{tabular}
    \caption{List of all control parameters, along with their values, used in the proposed flight controller}
    \label{control params: table}
\end{table}
These simulation results are for the following three cases:
\begin{enumerate}
    \item \textit{Case 1:} When the target is stationary at position $\,^\mathcal{I}\boldsymbol{r}_{\text{t}} = [0.0, 0.0, 0.0]^\top$ m in the inertial frame $\mathcal{I}$.
    \item \textit{Case 2:} When the target is moving with a constant linear velocity $\,^\mathcal{I}\boldsymbol{v}_{\text{t}} = [0.5, -0.5, 0.1]^\top$ m/s.
    \item \textit{Case 3:} When the target is undergoing circular motion on a horizontal plane with a turning radius of $1$ m and turning speed of $0.1$ rad/s.
\end{enumerate}

\subsection{Simulation without Uncertainty and Perturbation}
The performance of the proposed flight controller is first assessed under idealized conditions, with varying initial states and no model uncertainties or external disturbances.

\subsubsection{When the target is stationary}
\begin{figure*}[hbt!]
\centering
    \includegraphics[width=1.0 \textwidth]{Figures/01122025/Stationary/Linear_state_wo_d.png}
    \caption{This graph shows the evolution of controllable UAV pose $\,^\mathcal{I} \boldsymbol{q} = [\,^\mathcal{I} x, \,^\mathcal{I} y, \,^\mathcal{I} z, \psi]^\top$, along with its first derivative $\,^\mathcal{I} \dot{\boldsymbol{q}} = [\,^\mathcal{I} v_x, \,^\mathcal{I} v_y, \,^\mathcal{I} v_z, \dot{\psi}]^\top$, when the target is stationary. It also shows output of the outer loop controller, $\,^\mathcal{I} \boldsymbol{a}_\text{d} = [\,^\mathcal{I} {a}_{x\text{d}}, \,^\mathcal{I} {a}_{y\text{d}}, \,^\mathcal{I} {a}_{z\text{d}}]^\top$ and $[\phi_{\text{d}}, \theta_{\text{d}}, T_u]^\top$} 
    \label{Stationary_linear_wo_d: figure}
\end{figure*}

Fig. \ref{Stationary_linear_wo_d: figure} shows that the UAV has successfully maneuvered from the initial position $\,^\mathcal{I} \boldsymbol{r}_{\text{b}}(0) = [1.0, -1.0, 5.0]^\top$ m and the initial heading $\psi(0) = 10.0^\circ$ to the stationary target position $\,^\mathcal{I}\boldsymbol{r}_{\text{t}} = [0.0, 0.0, 0.0]^\top$ m and the target heading $\psi_\text{t} = 0.0^\circ$, hence demonstrating precise target tracking and orientation alignment using the proposed flight controller.
It also shows that the UAV linear velocity $\,^\mathcal{I} \boldsymbol{v} = [\,^\mathcal{I} v_x, \,^\mathcal{I} v_y, \,^\mathcal{I} v_z]^\top$ and yaw rate $\dot{\psi}$ simultaneously converges to the origin, hence demonstrating soft landing of the UAV using the proposed flight controller.

\begin{figure*}[hbt!]
\centering
    \includegraphics[width=1.0 \textwidth]{Figures/01122025/Stationary/Image_state_wo_d.png}
    \caption{This graph shows the evolution of target normalized position $\hat{\boldsymbol{r}} = [\hat{x}, \hat{y}]^{\top}$ and optical flow $\boldsymbol{h} = [h_{x}, h_{y}, h_{r}]^{\top}$, along with the evolution of their respective performance-constrained errors, when the target is stationary. 
    The last row shows evolution of target orientation $\alpha$, along with its first derivative since $\dot{\alpha} = \dot{\psi}$.}
    \label{Stationary_image_wo_d: figure}
\end{figure*}

Fig. \ref{Stationary_image_wo_d: figure} shows the performance constrained target normalized position $\hat{\boldsymbol{r}} = [\hat{x}, \hat{y}]^{\top}$ and optical flow $\boldsymbol{h} = [h_{x}, h_{y}, h_{r}]^{\top}$.
The convergence of $\hat{\boldsymbol{r}}$ to $[0.0, 0.0]^{\top}$ showcases precise target tracking, while simultaneous convergence of $\boldsymbol{h}$ to $[0.0, 0.0, -0.3]^{\top}$ demonstrate soft precise landing using these scale-independent target parameters as control feedback to the outer-loop controller.
Moreover, convergence of the target orientation $\alpha$ to the desired value $\alpha_\text{d}$ showcases UAV orientation alignment with that of the target without explicit knowledge of the target heading $\psi_t$.

\begin{figure*}[hbt!]
\centering
    \includegraphics[width=1.0 \textwidth]{Figures/01122025/Stationary/3d_wo_d.png}
    \caption{This 3-D plot shows the evolution of evolution of UAV position in inertial frame $\,^\mathcal{I} \boldsymbol{r} = [\,^\mathcal{I} x, \,^\mathcal{I} y, \,^\mathcal{I} z]^\top$, when the target is stationary.} 
    \label{Stationary_3d_wo_d: figure}
\end{figure*}

Fig. \ref{Stationary_3d_wo_d: figure} shows a 3-D plot depicting UAV motion from initial position $\,^\mathcal{I} \boldsymbol{r}_{\text{b}}(0) = [1.0, -1.0, 5.0]^\top$ m to the stationary target position $\,^\mathcal{I}\boldsymbol{r}_{\text{t}} = [0.0, 0.0, 0.0]^\top$ m.

\subsubsection{ When the target is moving with a constant linear velocity}
\begin{figure*}[hbt!]
\centering
    \includegraphics[width=1.0 \textwidth]{Figures/01122025/Linear_positive/Linear_state_wo_d.png}
    \caption{This graph shows the evolution of controllable UAV pose $\,^\mathcal{I} \boldsymbol{q} = [\,^\mathcal{I} x, \,^\mathcal{I} y, \,^\mathcal{I} z, \psi]^\top$, along with its first derivative $\,^\mathcal{I} \dot{\boldsymbol{q}} = [\,^\mathcal{I} v_x, \,^\mathcal{I} v_y, \,^\mathcal{I} v_z, \dot{\psi}]^\top$, when the target is moving linearly with constant velocity. It also shows output of the outer loop controller, $\,^\mathcal{I} \boldsymbol{a}_\text{d} = [\,^\mathcal{I} {a}_{x\text{d}}, \,^\mathcal{I} {a}_{y\text{d}}, \,^\mathcal{I} {a}_{z\text{d}}]^\top$ and $[\phi_{\text{d}}, \theta_{\text{d}}, T_u]^\top$} 
    \label{Linear_linear_wo_d: figure}
\end{figure*}

Fig. \ref{Linear_linear_wo_d: figure} shows that the UAV has successfully maneuvered from initial position $\,^\mathcal{I} \boldsymbol{r}_{\text{b}}(0) = [1.0, -1.0, 5.0]^\top$ m to the target moving with a constant linear velocity $\,^\mathcal{I}\boldsymbol{v}_{\text{t}} = [0.5, -0.5, 0.1]^\top$ m/s. This graph demonstrates precise target tracking capability of the proposed flight controller, even when the target is not stationary.
It also shows that the UAV linear velocity $\,^\mathcal{I} \boldsymbol{v}$ successfully converges to the target velocity $\,^\mathcal{I}\boldsymbol{v}_{\text{t}}$, hence demonstrating soft landing of the UAV using the proposed flight controller when the target is not stationary.

\begin{figure*}[hbt!]
\centering
    \includegraphics[width=1.0 \textwidth]{Figures/01122025/Linear_positive/Image_state_wo_d.png}
    \caption{This graph shows the evolution of target normalized position $\hat{\boldsymbol{r}} = [\hat{x}, \hat{y}]^{\top}$ and optical flow $\boldsymbol{h} = [h_{x}, h_{y}, h_{r}]^{\top}$, along with the evolution of their respective performance-constrained errors, when the target is moving linearly with constant velocity. 
    The last row shows evolution of target orientation $\alpha$, along with its first derivative since $\dot{\alpha} = \dot{\psi}$.}
    \label{Linear_image_wo_d: figure}
\end{figure*}

Fig. \ref{Linear_image_wo_d: figure} shows the performance constrained target normalized position $\hat{\boldsymbol{r}}$ and optical flow $\boldsymbol{h}$.
Similar to Fig. \ref{Stationary_image_wo_d: figure}, the scale-independent target parameters $\hat{\boldsymbol{r}}$ and $\boldsymbol{h}$ converge to $[0.0, 0.0]^{\top}$ and $[0.0, 0.0, -0.3]^{\top}$, respectively. This convergence demonstrates successful soft precise landing of the UAV on the mobile platform using the proposed flight controller.

\begin{figure*}[hbt!]
\centering
    \includegraphics[width=1.0 \textwidth]{Figures/01122025/Linear_positive/3d_wo_d.png}
    \caption{This 3-D plot shows the evolution of evolution of UAV position in inertial frame $\,^\mathcal{I} \boldsymbol{r} = [\,^\mathcal{I} x, \,^\mathcal{I} y, \,^\mathcal{I} z]^\top$, when the target is moving linearly with constant velocity.} 
    \label{Linear_3d_wo_d: figure}
\end{figure*}

Fig. \ref{Linear_3d_wo_d: figure} shows a 3-D plot depicting UAV motion from initial position $\,^\mathcal{I} \boldsymbol{r}_{\text{b}}(0) = [1.0, -1.0, 5.0]^\top$ m to the target moving with a constant linear velocity $\,^\mathcal{I}\boldsymbol{v}_{\text{t}} = [0.5, -0.5, 0.1]^\top$ m/s.

\subsubsection{When the target is undergoing circular motion on a horizontal plane}
\begin{figure*}[hbt!]
\centering
    \includegraphics[width=1.0 \textwidth]{Figures/01122025/Circular/Linear_state_wo_d.png}
    \caption{This graph shows the evolution of controllable UAV pose $\,^\mathcal{I} \boldsymbol{q} = [\,^\mathcal{I} x, \,^\mathcal{I} y, \,^\mathcal{I} z, \psi]^\top$, along with its first derivative $\,^\mathcal{I} \dot{\boldsymbol{q}} = [\,^\mathcal{I} v_x, \,^\mathcal{I} v_y, \,^\mathcal{I} v_z, \dot{\psi}]^\top$, when the target is undergoing circular motion on a horizontal plane. It also shows output of the outer loop controller, $\,^\mathcal{I} \boldsymbol{a}_\text{d} = [\,^\mathcal{I} {a}_{x\text{d}}, \,^\mathcal{I} {a}_{y\text{d}}, \,^\mathcal{I} {a}_{z\text{d}}]^\top$ and $[\phi_{\text{d}}, \theta_{\text{d}}, T_u]^\top$} 
    \label{Circular_linear_wo_d: figure}
\end{figure*}

\begin{figure*}[hbt!]
\centering
    \includegraphics[width=1.0 \textwidth]{Figures/01122025/Circular/Image_state_wo_d.png}
    \caption{This graph shows the evolution of target normalized position $\hat{\boldsymbol{r}} = [\hat{x}, \hat{y}]^{\top}$ and optical flow $\boldsymbol{h} = [h_{x}, h_{y}, h_{r}]^{\top}$, along with the evolution of their respective performance-constrained errors, when the target is undergoing circular motion on a horizontal plane. 
    The last row shows evolution of target orientation $\alpha$, along with its first derivative since $\dot{\alpha} = \dot{\psi}$.}
    \label{Circular_image_wo_d: figure}
\end{figure*}

\begin{figure*}[hbt!]
\centering
    \includegraphics[width=1.0 \textwidth]{Figures/01122025/Circular/3d_wo_d.png}
    \caption{This 3-D plot shows the evolution of evolution of UAV position in inertial frame $\,^\mathcal{I} \boldsymbol{r} = [\,^\mathcal{I} x, \,^\mathcal{I} y, \,^\mathcal{I} z]^\top$, when the target is undergoing circular motion on a horizontal plane.} 
    \label{Circular_3d_wo_d: figure}
\end{figure*}

Similar to the previous two cases, Figs. \ref{Circular_linear_wo_d: figure} - \ref{Circular_3d_wo_d: figure} shows that the UAV can achieve soft precise landing using the proposed flight controller when the target is undergoing circular motion on a horizontal plane.

\subsection{Simulation with Uncertainty and Perturbation}
In order to simulate more realistic scenarios, the following real-world settings are added in the numerical simulation is conducted in MATLAB under following scenarios:
\begin{enumerate}
    \item \underline{Measurement Noise}: The estimated optical flow $\boldsymbol{h}$ is corrupted using additive Gaussian white noise with $\mu$ = 0.1 rad/s and $\sigma$ = 0.05 rad/s.
    \item \underline{Ground Effect}: The ground effect factor is added to the the total rotor control thrust $\,^\mathcal{B} T_u$ as given below:
    \begin{equation} \label{ground effect: equation}
        \,^\mathcal{B} T_{\infty} = \frac{1}{1- \left(\frac{R}{4\,^\mathcal{I} z_{\text{t}}} \right)^2} \,^\mathcal{B} T_u, \forall~t \geq 0,
    \end{equation}
    where  $R$ is the propeller radius.
    \item \underline{Computational Delay}: A delay of 1 time step is added during the computation of the combined thrust-torque command $\boldsymbol{u} = [\,^\mathcal{B} {T}_u, \,^\mathcal{B} \boldsymbol{\tau}_u^\top]^\top$.
    \item \underline{Visibility Constraint}: The target image error $\hat{\boldsymbol{r}}_\text{e}$ is constrained using the performance function $\boldsymbol{p}_1(t)$ as shown in \eqref{prescribed perfromance: equation}.
    \item \underline{Input Constraint}: Given the motor-mixing matrix $M$, the motor thrust vector $\boldsymbol{T}(t)$ is computed as:
    \begin{equation}
        \boldsymbol{T}(t) = M^{-1}\boldsymbol{u},
    \end{equation}
    where $T_{\min} \leq T_k(t) \leq T_{\max}$, for motor $k = \{1,2,3,4\}$, with $T_{\min} = 0.01 Nm$ and $T_{\max} = 7.0 Nm$.
\end{enumerate}

\subsubsection{When the target is stationary}
\begin{figure*}[hbt!]
\centering
    \includegraphics[width=1.0 \textwidth]{Figures/01122025/Stationary/Linear_state_w_d.png}
    \caption{This graph shows the evolution of controllable UAV pose $\,^\mathcal{I} \boldsymbol{q} = [\,^\mathcal{I} x, \,^\mathcal{I} y, \,^\mathcal{I} z, \psi]^\top$, along with its first derivative $\,^\mathcal{I} \dot{\boldsymbol{q}} = [\,^\mathcal{I} v_x, \,^\mathcal{I} v_y, \,^\mathcal{I} v_z, \dot{\psi}]^\top$, when the target is stationary. It also shows output of the outer loop controller, $\,^\mathcal{I} \boldsymbol{a}_\text{d} = [\,^\mathcal{I} {a}_{x\text{d}}, \,^\mathcal{I} {a}_{y\text{d}}, \,^\mathcal{I} {a}_{z\text{d}}]^\top$ and $[\phi_{\text{d}}, \theta_{\text{d}}, T_u]^\top$} 
    \label{Stationary_linear_w_d: figure}
\end{figure*}

Fig. \ref{Stationary_linear_w_d: figure} shows that fluctuation in the desired acceleration $\,^\mathcal{I}\boldsymbol{a}_{\text{d}}$ plot, depicting the impact of introducing model uncertainty and exogenous perturbation in the system dynamics. In spite of these fluctuations, the UAV has successfully maneuvered from the initial position $\,^\mathcal{I} \boldsymbol{r}_{\text{b}}(0) = [1.0, -1.0, 5.0]^\top$ m and the initial heading $\psi(0) = 10.0^\circ$ to the stationary target position $\,^\mathcal{I}\boldsymbol{r}_{\text{t}} = [0.0, 0.0, 0.0]^\top$ m and the target heading $\psi_\text{t} = 0.0^\circ$. This demonstrates robust performance of the proposed flight controller in precise target tracking and orientation alignment.
Similarly, the UAV linear velocity $\,^\mathcal{I} \boldsymbol{v} = [\,^\mathcal{I} v_x, \,^\mathcal{I} v_y, \,^\mathcal{I} v_z]^\top$ and yaw rate $\dot{\psi}$ simultaneously converges to the origin, hence demonstrating robust performance of the proposed flight controller in achieving soft landing.

\begin{figure*}[hbt!]
\centering
    \includegraphics[width=1.0 \textwidth]{Figures/01122025/Stationary/Image_state_w_d.png}
    \caption{This graph shows the evolution of target normalized position $\hat{\boldsymbol{r}} = [\hat{x}, \hat{y}]^{\top}$ and optical flow $\boldsymbol{h} = [h_{x}, h_{y}, h_{r}]^{\top}$, along with the evolution of their respective performance-constrained errors, when the target is stationary. 
    The last row shows evolution of target orientation $\alpha$, along with its first derivative since $\dot{\alpha} = \dot{\psi}$.}
    \label{Stationary_image_w_d: figure}
\end{figure*}

Fig. \ref{Stationary_image_w_d: figure} validates that the scale-independent target parameters  $\hat{\boldsymbol{r}}$ and $\boldsymbol{h}$ satisfy the prescribed-performance constraints, even in the presence of model uncertainty and exogenous perturbation in the system dynamics.

\begin{figure*}[hbt!]
\centering
    \includegraphics[width=1.0 \textwidth]{Figures/01122025/Stationary/3d_w_d.png}
    \caption{This 3-D plot shows the evolution of evolution of UAV position in inertial frame $\,^\mathcal{I} \boldsymbol{r} = [\,^\mathcal{I} x, \,^\mathcal{I} y, \,^\mathcal{I} z]^\top$, when the target is stationary.} 
    \label{Stationary_3d_w_d: figure}
\end{figure*}

Fig. \ref{Stationary_3d_wo_d: figure} shows a 3-D plot depicting smooth UAV motion from initial position $\,^\mathcal{I} \boldsymbol{r}_{\text{b}}(0) = [1.0, -1.0, 5.0]^\top$ m to the stationary target position $\,^\mathcal{I}\boldsymbol{r}_{\text{t}} = [0.0, 0.0, 0.0]^\top$ m, even in the presence of model uncertainty and exogenous perturbation in the system dynamics.

\subsubsection{ When the target is moving linearly with constant velocity}
\begin{figure*}[hbt!]
\centering
    \includegraphics[width=1.0 \textwidth]{Figures/01122025/Linear_positive/Linear_state_w_d.png}
    \caption{This graph shows the evolution of controllable UAV pose $\,^\mathcal{I} \boldsymbol{q} = [\,^\mathcal{I} x, \,^\mathcal{I} y, \,^\mathcal{I} z, \psi]^\top$, along with its first derivative $\,^\mathcal{I} \dot{\boldsymbol{q}} = [\,^\mathcal{I} v_x, \,^\mathcal{I} v_y, \,^\mathcal{I} v_z, \dot{\psi}]^\top$, when the target is moving linearly with constant velocity. It also shows output of the outer loop controller, $\,^\mathcal{I} \boldsymbol{a}_\text{d} = [\,^\mathcal{I} {a}_{x\text{d}}, \,^\mathcal{I} {a}_{y\text{d}}, \,^\mathcal{I} {a}_{z\text{d}}]^\top$ and $[\phi_{\text{d}}, \theta_{\text{d}}, T_u]^\top$} 
    \label{Linear_linear_w_d: figure}
\end{figure*}

\begin{figure*}[hbt!]
\centering
    \includegraphics[width=1.0 \textwidth]{Figures/01122025/Linear_positive/Image_state_w_d.png}
    \caption{This graph shows the evolution of target normalized position $\hat{\boldsymbol{r}} = [\hat{x}, \hat{y}]^{\top}$ and optical flow $\boldsymbol{h} = [h_{x}, h_{y}, h_{r}]^{\top}$, along with the evolution of their respective performance-constrained errors, when the target is moving linearly with constant velocity. 
    The last row shows evolution of target orientation $\alpha$, along with its first derivative since $\dot{\alpha} = \dot{\psi}$.}
    \label{Linear_image_w_d: figure}
\end{figure*}

\begin{figure*}[hbt!]
\centering
    \includegraphics[width=1.0 \textwidth]{Figures/01122025/Linear_positive/3d_w_d.png}
    \caption{This 3-D plot shows the evolution of evolution of UAV position in inertial frame $\,^\mathcal{I} \boldsymbol{r} = [\,^\mathcal{I} x, \,^\mathcal{I} y, \,^\mathcal{I} z]^\top$, when the target is moving linearly with constant velocity.} 
    \label{Linear_3d_w_d: figure}
\end{figure*}

The simulation results shown in Figs. \ref{Linear_linear_w_d: figure}-\ref{Linear_3d_w_d: figure} validate that the proposed flight controller is robust against model uncertainty and exogenous perturbation in achieving soft precise landing when the target has linear motion.

\subsubsection{When the target is undergoing circular motion on a horizontal plane}

\begin{figure*}[hbt!]
\centering
    \includegraphics[width=1.0 \textwidth]{Figures/01122025/Circular/Linear_state_w_d.png}
    \caption{This graph shows the evolution of controllable UAV pose $\,^\mathcal{I} \boldsymbol{q} = [\,^\mathcal{I} x, \,^\mathcal{I} y, \,^\mathcal{I} z, \psi]^\top$, along with its first derivative $\,^\mathcal{I} \dot{\boldsymbol{q}} = [\,^\mathcal{I} v_x, \,^\mathcal{I} v_y, \,^\mathcal{I} v_z, \dot{\psi}]^\top$, when the target is undergoing circular motion on a horizontal plane. It also shows output of the outer loop controller, $\,^\mathcal{I} \boldsymbol{a}_\text{d} = [\,^\mathcal{I} {a}_{x\text{d}}, \,^\mathcal{I} {a}_{y\text{d}}, \,^\mathcal{I} {a}_{z\text{d}}]^\top$ and $[\phi_{\text{d}}, \theta_{\text{d}}, T_u]^\top$} 
    \label{Circular_linear_w_d: figure}
\end{figure*}

\begin{figure*}[hbt!]
\centering
    \includegraphics[width=1.0 \textwidth]{Figures/01122025/Circular/Image_state_w_d.png}
    \caption{This graph shows the evolution of target normalized position $\hat{\boldsymbol{r}} = [\hat{x}, \hat{y}]^{\top}$ and optical flow $\boldsymbol{h} = [h_{x}, h_{y}, h_{r}]^{\top}$, along with the evolution of their respective performance-constrained errors, when the target is undergoing circular motion on a horizontal plane. 
    The last row shows evolution of target orientation $\alpha$, along with its first derivative since $\dot{\alpha} = \dot{\psi}$.}
    \label{Circular_image_w_d: figure}
\end{figure*}

\begin{figure*}[hbt!]
\centering
    \includegraphics[width=1.0 \textwidth]{Figures/01122025/Circular/3d_w_d.png}
    \caption{This 3-D plot shows the evolution of evolution of UAV position in inertial frame $\,^\mathcal{I} \boldsymbol{r} = [\,^\mathcal{I} x, \,^\mathcal{I} y, \,^\mathcal{I} z]^\top$, when the target is undergoing circular motion on a horizontal plane.} 
    \label{Circular_3d_w_d: figure}
\end{figure*}

The simulation results shown in Figs. \ref{Circular_linear_w_d: figure}-\ref{Circular_3d_w_d: figure} validate that the proposed flight controller is robust against model uncertainty and exogenous perturbation in achieving soft precise landing when the target has circular motion.

% \begin{thebibliography}{99}

% % \bibitem{Keshavan} 
% % J.~Keshavan, ``Adaptive Visual Control Strategies for Autonomous Landing of Quadrotor Platforms,".

% \bibitem{Lin} 
% J.~Lin, Y. Wang, Z. Miao, R. Fierro, "Low-complexity control for vision-based landing of Quadrotor UAV on unknown platform."

% \end{thebibliography}

\end{document}